\subsection{区间的刻画}

引理 6. --- 设 D 是一个可数的全序集，不退化为单元素集，具有最小元和最大元。假设 D 无间隙（E, III, p.~73, 习题 19），即，每个开区间 {]}x, y{[}，其中 x 和 y 是 D 中满足 x \textless{} y 的元素，都非空。于是存在一个从 I ∩ Q 到 D 的有序集同构。

令 \(a\) 为最小元，\(b\) 为最大元。依假设，有 \(b \neq a\)，且 \(]a\) , \(x[\neq \emptyset\) 对每个 \(x \in \mathrm{D} - \{a\}\) 成立。集合 \(\mathrm{D} - \{a\}\) 是全序的、非空且没有最小元；因此它是无限的（E, III, p.~34, 命题 3 的推论 1）。集合 \(\mathbf{I} \cap \mathbf{Q}\) 和 D 都是无限可数的，且与 \(\mathbf{N}\) 等势。

选取双射 \(n \mapsto a_{n}\) 与 \(n \mapsto b_{n}\)，分别从 N 到 \(I \cap Q\) 和 D，使得 \(a_{0} = 0\)、\(a_{1} = 1\)、\(b_{0} = a\)、\(b_{1} = b\)。存在唯一一个严格递增映射 \(f: I \cap Q \to D\)，具有以下性质：\(f(0) = a\) 且 \(f(1) = b\)；当 \(n \geqslant 2\) 时，\(f(a_{n}) = b_{m}\)，其中 m 是使得从 \(\{a_{0}, \ldots, a_{n}\}\) 到 D 的、在 \(\{a_{0}, \ldots, a_{n-1}\}\) 上与 f 一致且将 \(a_{n}\) 映为 \(b_{m}\) 的映射严格递增的最小自然整数。这些性质确实通过对 n 归纳定义了 \(f(a_{n})\)，而整数 m 的存在性由 D 无间隙这一事实保证。

由于映射 \(f\) 严格递增且 \(\mathbf{I} \cap \mathbf{Q}\) 是全序的，\(f\) 定义了从 \(\mathbf{I} \cap \mathbf{Q}\) 到其像的有序集同构（E, III, p.~14, 命题 11）。还需证明 \(f\) 是满射。为此，我们证明 \(b_{m}\) 属于 \(f\) 的像，对所有 \(m \in \mathbf{N}\) 成立。

我们有 \(b_{0}=f(0)\) 和 \(b_{1}=f(1)\)。假设 \(m\geqslant2\)，并且对于 \(0\leqslant k\leqslant m-1\)，存在 \(c_{k}\in I\cap Q\) 使得 \(f(c_{k})=b_{k}\)。由于 f 是单射，有 \(c_{0}=0\) 和 \(c_{1}=1\)。考虑满足下列条件的最小整数 \(n\in N\)：\(a_{n}\) 不属于 \(\{c_{0},\ldots,c_{m-1}\}\)，并且从 \(\{c_{0},\ldots,c_{m-1},a_{n}\}\) 到 D 的、在 \(\{c_{0},\ldots,c_{m-1}\}\) 上与 f 一致且将 \(a_{n}\) 映为 \(b_{m}\) 的映射 g 严格递增；这样的整数存在，因为 \(I\cap Q\) 无间隙。令 \(f'\) 为从 \(\{a_{0},\ldots,a_{n}\}\) 到 D 的映射，它在 \(\{a_{0},\ldots,a_{n-1}\}\) 上与 f 一致，并将 \(a_{n}\) 映为 \(b_{m}\)。我们来证明它严格递增。令 \(j\in\{0,\ldots,n-1\}\)；由整数 n 的定义，存在 \(k\in\{0,\ldots,m-1\}\)，使得 \(a_{j}=c_{k}\)，或者 \(a_{j}<c_{k}\) 且 \(b_{m}\geqslant f(c_{k})\)，或者 \(a_{j}>c_{k}\) 且 \(b_{m}\leqslant f(c_{k})\)。假设 \(b_{k}<b_{m}\)；于是 \(g(c_{k})=f(c_{k})=b_{k}<b_{m}=g(a_{n})\)，从而由于 g 严格递增，有 \(c_{k}<a_{n}\)，进而 \(a_{j}\leqslant c_{k}<a_{n}\)；此外，\(f'(a_{j})=f(a_{j})\leqslant f(c_{k})=b_{k}<b_{m}=f'(a_{n})\)。同样，若 \(b_{k}>b_{m}\)，则有 \(a_{n}<c_{k}\leqslant a_{j}\)，且 \(f'(a_{j})=f(a_{j})\geqslant f(c_{k})=b_{k}>b_{m}=f'(a_{n})\)。由于 f 本身严格递增，映射 \(f'\) 也严格递增。

若 \(m'\) 是满足 \(f(a_n) = b_{m'}\) 的整数，则根据 f 的定义，有 \(m' \leqslant m\)。若有 \(m' < m\)，则有 \(f(a_n) = b_{m'} = f(c_{m'})\)，从而由于 f 是单射，得到 \(a_n = c_{m'}\)，这与 \(a_n\) 的定义矛盾。因此 \(m' = m\) 且 \(f(a_n) = b_m\)，这证明了 \(b_m\) 属于 f 的像，并完成了对 f 为满射的归纳证明。

命题 16. --- 设 E 是一个不退化为单元素集的全序集。假设 E 的每个子集都有最小上界，并且存在一个可数子集与形如 {]}x, y{[} 的每个开区间相交，其中 x 和 y 是 E 中满足 x \textless{} y 的元素。于是存在从 I 到 E 的有序集同构。

由于 ∅ 和 E 在 E 中各自都有最小上界，E 具有最小元 a 和最大元 b。它们彼此不同，因为 E 不退化为单元素集。设 D' 是 E 的一个可数子集，与 E 中形如 {]}x, y{[} 且满足 x \textless{} y 的每个开区间相交。置 D = D' ∪ \{a, b\}。集合 D 是全序的且无间隙。由引理 6，存在从 I ∩ Q 到 D 的序集同构 f。

对每个 \(t \in I\)，令 \(g(t)\) 为 \(f([0, t] \cap \mathbf{Q})\) 在 E 中的上界。对每个 \(x \in E\)，令 \(h(x)\) 为 \(f^{-1}([a, x] \cap D)\) 在 I 中的上界。如此定义的映射 \(g: I \to E\) 和 \(h: E \to I\) 都是递增的，g 在 \(I \cap Q\) 上与 f 一致，而 h 在 D 上与 \(f^{-1}\) 一致。

因此有 \(g(h(y)) = y\)，对所有 \(y \in D\) 成立。令 \(x \in E\)。若有 \(g(h(x)) > x\)，则开区间 {]}x, \(g(h(x))[\) 将包含 D 中的某个点 y，而关系 \(g(h(y)) = y < g(h(x))\) 将与 \(g \circ h\) 递增这一事实矛盾。同样，不可能有 \(g(h(x)) < x\)。因此 \(g(h(x)) = x\)，这证明了 \(g \circ h\) 是 E 的恒等映射。类似地可证明 \(h \circ g\) 是 I 的恒等映射。因此，\(g: I \to E\) 与 \(h: E \to I\) 是互为逆的有序集同构。

注. --- 设 E 是一个全序集。E 的开区间集合（有界或无界）对有限交稳定。它是 E 上拓扑 \(\mathcal{T}_{0}(\mathrm{E})\) 的一个基（TG, I, p.~91, 习题 5）。拓扑 \(\mathcal{T}_{0}(\mathbf{I})\) 与 \(\mathbf{I}\) 上由 \(\mathbf{R}\) 的拓扑诱导的拓扑相同。因此，在命题 16 中，从 \(\mathbf{I}\) 到 E 的每个有序集同构，都是从 \(\mathbf{I}\) 到赋予 E 以拓扑 \(\mathcal{T}_{0}(\mathrm{E})\) 所得拓扑空间的同胚。

推论. --- 设 R 是 I 上的一个等价关系。下列条件等价：

\begin{enumerate}
\def\labelenumi{(\roman{enumi})}
\item
  R 的每个等价类都是 I 的一个闭区间，且不同于 I；
\item
  存在一个递增满射 \(u\colon\mathbf{I}\to\mathbf{I}\)，使得 R 是与 \(u\) 相伴的等价关系。
\end{enumerate}

这样的映射 u 一旦存在就是连续的，并且通过取商诱导出从 I/R 到 I 的同胚。

记 \(p: I \rightarrow I/R\) 为典范满射。

假设条件 (i) 成立。对于 R 的等价类 A 和 B，记 A \textless{} B，当有 a \textless{} b 对所有 \(a \in A\) 及每个 \(b \in B\) 时。在 I/R 中，关系“ A = B 或 A \textless{} B”是一种序关系。事实上，它是自反的；它是反对称的，因为不可能同时有 A \textless{} B 和 B \textless{} A；它是传递的，因为关系 A \textless{} B 和 B \textless{} C 蕴含 A \textless{} C。若 A 和 B 是 I/R 中彼此不同的元素，则它们是 I 的不交闭区间，因而 A \textless{} B 或 B \textless{} A。赋予 I/R 以上述定义的序关系后，I/R 因而是全序的。映射 \(p: I \to I/R\) 是递增的。

由于 (i)，集合 I/R 不退化为单元素集。

设 F 是 I/R 的一个子集。我们来证明 F 在 I/R 中有最小上界。置 \(F' = \overline{p}^{-1}(F)\)；记 a 为 \(F'\) 在 I 中的最小上界，记 A 为 a 关于 R 的等价类。由于 a 是 \(F'\) 在 I 中的上界，A 是 F 在 I/R 中的上界。反之，设 \(B \in I/R\) 是 F 的一个上界；置 \(b = \sup(B)\)。于是 \(F'\) 的每个元素都被 b 上界。因而有 \(a \leqslant b\)。由于 R 下的等价类是闭区间，a 属于 A，b 属于 B。因此有 \(A = p(a) \leqslant p(b) = B\)。这证明了 A 是 F 的最小上界。

设 A 和 B 是 I/R 中满足 A \textless{} B 的元素。令 a 为 A 的最小上界，令 b 为 B 的最大下界。由于 \(a \in A\) 且 \(b \in B\)，有 a \textless{} b。{]}a, b{[} 中任一元素关于 R 的等价类，都是 I/R 的区间 {]}A, B{[} 中的元素。由于 \(I \cap Q\) 与 {]}a, b{[} 相交，其在 p 下的像与 {]}A, B{[} 相交。

这样，我们证明了全序集 I/R 满足命题 16 的假设。因此存在一个有序集同构 \(f \colon \mathbf{I} / \mathrm{R} \to \mathbf{I}\)。映射 \(u = f \circ p\) 是从 I 到 I 的满射递增映射，而与 u 相伴的等价关系就是关系 R；这证明了条件 (ii) 成立。

反之，假设条件 (ii) 成立。设 \(u: I \rightarrow I\) 是一个递增满射，且 R 是与 u 相伴的等价关系。

令 \(a\) 为 \(\mathbf{I}\) 中的一点。集合 \(\mathrm{A} = \overline{u}^{-1}(a)\) 是 \(\mathbf{I}\) 的一个区间，因为映射 \(u\) 是递增的。由于 \(u\) 是满射，A 既非空也不等于 \(\mathbf{I}\)。令 \(b\) 为 A 在 \(\mathbf{I}\) 中的最小上界。有 \(u(b) \geqslant a\)。若 \(u(b) > a\)，存在 \(c \in ]a, u(b)[\) 及 \(d \in \mathbf{I}\)，使得 \(u(d) = c\)，因为 \(u\) 是满射。由于 \(u\) 是递增的且 \(a < u(d) < u(b)\)，\(d\) 是 A 的每个元素的上界，且有 \(d < b\)，这与 \(b\) 是 A 的上确界这一假设矛盾。因此 \(u(b) = a\)，即 \(b \in \mathrm{A}\)。类似地可证明 A 包含其下界。因此，集合 A 是 \(\mathbf{I}\) 的一个不同于 \(\mathbf{I}\) 的闭区间。这证明了条件 (i) 成立。

现在证明映射 \(u\) 是连续的。只需证明，对每个 \(a \in \mathbf{I}\)，集合 \(\overline{u}^{-1}([a, \to[)\) 和 \(\overline{u}^{-1}([\leftarrow, a[)\) 都是开集。令 \(b\) 为 \(\overline{u}^{-1}(a)\) 的上界；有 \(u(b) = a\)。若 \(x \in \mathbf{I}\) 满足 \(u(x) > a\)，则必有 \(x > b\)；反之，若 \(x > b\)，则有 \(u(x) \geqslant a\)，并且 \(u(x) \neq a\)，因为 \(b\) 是 \(\overline{u}^{-1}(a)\) 的最小上界。因此有 \(\overline{u}^{-1}([a, \to[) = ]b, \to[\)，这表明 \(u\) 对区间 \(]a, \to[\) 的逆像是开集。类似地可证明 \(]\leftarrow, a[\) 的逆像是开集。由此映射 \(u\) 连续。

映射 \(v \colon \mathbf{I} / \mathrm{R} \to \mathbf{I}\) 是由 \(u\) 取商得到的，于是是连续双射。由于 \(\mathbf{I}\) 是紧的，\(\mathbf{I} / \mathrm{R}\) 是拟紧的（TG, I, p.~62, 定理 2），且 \(v\) 是同胚（TG, I, p.~63, 推论 2）。

命题 17. --- 设 X 是一个连通紧拓扑空间。设 a 是 X 中的一点，U 是 \(X - \{a\}\) 的一个非空开闭子集。

\begin{enumerate}
\def\labelenumi{\alph{enumi})}
\item
  U 在 X 中的闭包 \(\overline{U}\) 等于 \(U \cup \{a\}\)，并且是连通的。
\item
  设 \(a'\) 是 X 中异于 \(a\) 的一点，\(U'\) 是 \(X - \{a'\}\) 的一个非空开闭子集。若 \(a \notin U'\) 且 \(\overline{U} \cap \overline{U'} \neq \varnothing\)，则有 \(\overline{U'} \subset U\)。反之，若 \(a \in U'\) 且 \(X \neq U \cup U'\)，则有 \(\overline{U} \subset U'\)。
\item
  U 中存在一点 \(b\)，使得 \(X - \{b\}\) 连通。
\end{enumerate}

我们来证明断言 \(a\)。令 V 为 U 在 \(X - \{a\}\) 中的补集。由于 U 在 \(X - \{a\}\) 中是闭的，V 在 \(X - \{a\}\) 中是开的，从而在 X 中也是开的。因此有 \(U \subset \overline{\mathrm{U}} \subset X - V = U \cup \{a\}\)。同样，U 是 X 的开子集。由于 X 是连通的，U 在 X 中不是闭的，从而 \(\overline{\mathrm{U}} = U \cup \{a\}\)。类似地，有 \(\overline{\mathrm{V}} = V \cup \{a\}\)。

设 F 和 G 是 \(\overline{U}\) 的不交闭子集，且 \(\overline{U} = F \cup G\)。假设 \(a \in G\)，我们来证明 F 为空；若 \(a \in F\)，论证相同。集合 \(G \cup V = G \cup \overline{V}\) 是 X 的闭子集，并且与 F 不交；此外有 \(X = \overline{U} \cup V = F \cup (G \cup V)\)。由于 X 连通且 G 非空，F 为空。这证明了 \(\overline{U}\) 是连通的。

我们来证明 \(b\)。假设 \(a \notin \mathrm{U}'\) 且 \(\overline{\mathrm{U}} \cap \overline{\mathrm{U}'} \neq \emptyset\)。由断言 \(a\)，有 \(\overline{\mathrm{U}} = \mathrm{U} \cup \{a\}\) 和 \(\overline{\mathrm{U}'} = \mathrm{U}' \cup \{a'\}\)。由于 \(a \notin \mathrm{U}'\) 且 \(a \neq a'\)，集合 U 与 \(\overline{\mathrm{U}'}\) 有公共点。仍由 \(a\)，\(\overline{\mathrm{U}'}\) 是 X 的连通子集；由于它包含于 \(\mathrm{X} - \{a\}\) 且与开闭子集 U 相交，有 \(\overline{\mathrm{U}'} \subset \mathrm{U}\)，这正是所要证明的。第二个断言由第一个断言分别考虑 \(\overline{\mathrm{U}}\) 和 \(\overline{\mathrm{U}'}\) 在 X 中的补集而得。

最后来证明断言 \(c\)。反设对每个 \(x \in \mathrm{U}\)，集合 \(\mathrm{X} - \{x\}\) 都不连通，并选取开且闭、彼此不交且非空的子集 \(\mathrm{U}_x\) 和 \(\mathrm{V}_x\)，它们属于 \(\mathrm{X} - \{x\}\)，使得 \(\mathrm{X} - \{x\} = \mathrm{U}_x \cup \mathrm{V}_x\) 且 \(a \in \mathrm{V}_x\)。根据断言 \(b\)，将 \(\mathrm{U}\) 和 \(\mathrm{U}_x\) 分别看作 \(\mathrm{X} - \{a\}\) 和 \(\mathrm{X} - \{x\}\) 中的开闭子集，得到 \(\overline{\mathrm{U}_x} \subset \mathrm{U}\)，对所有 \(x \in \mathrm{U}\) 成立。设 \(x\) 和 \(y\) 是 U 中满足 \(x \in \mathrm{U}_y\) 的点；于是 \(x \neq y\)。仍根据断言 \(b\)，将 \(\mathrm{U}_x\) 和 \(\mathrm{U}_y\) 分别看作 \(\mathrm{X} - \{x\}\) 和 \(\mathrm{X} - \{y\}\) 中的开闭子集，关系 \(x \in \mathrm{U}_y\) 蕴含关系 \(\overline{\mathrm{U}_x} \subset \overline{\mathrm{U}_y}\)。因此，关系 \(x \in \mathrm{U}_y\) 与 \(\overline{\mathrm{U}_x} \subset \overline{\mathrm{U}_y}\) 等价。

由于空间 X 是紧的，U 的满足 \(\bigcap_{x\in\mathrm{S}}\overline{\mathrm{U}_{x}}\neq\varnothing\) 的子集 S 的集合具有有限特征（E, III, p.~34）。由 E, III, p.~35, 定理 1，存在 U 的一个极大子集 S，使得 \(\mathrm{C}=\bigcap_{x\in\mathrm{S}}\overline{\mathrm{U}_{x}}\) 非空。令 \(c\) 为 C 中的一点。对 S 的每个元素 \(x\)，都有 \(c\in\overline{\mathrm{U}_{x}}\)，从而 \(c\in\mathrm{U}\)；于是 \(\overline{\mathrm{U}_{c}}\subset\overline{\mathrm{U}_{x}}\)。因此有 \(\overline{\mathrm{U}_{c}}\subset\mathrm{C}\)；再由 S 的极大性，\(\mathrm{C}=\overline{\mathrm{U}_{c}}\)。对所有满足 \(x\in\mathrm{C}\) 且 \(x\neq c\) 的情形，还有 \(\overline{\mathrm{U}_{x}}\subset\overline{\mathrm{U}_{c}}\) 且 \(\overline{\mathrm{U}_{x}}\neq\overline{\mathrm{U}_{c}}\)。由 S 的极大性，\(\mathrm{C}=\{c\}\)，因而 \(\overline{\mathrm{U}_{c}}=\{c\}\)，这与 \(\mathrm{U}_{c}\) 是 \(\mathrm{X} - \{c\}\) 的非空开闭子集这一假设矛盾。

推论. --- 设 X 是一个紧连通拓扑空间，N 是补集连通的 X 中各点组成的集合。集合 X 是包含 N 的唯一紧连通子集。

设 S 是 X 的连通紧子集，满足 \(N \subset S\)。假设 \(S \neq X\)，并令 \(x \in X - S\)。根据假设，\(X - \{x\}\) 不连通；因此存在 \(X - \{x\}\) 中彼此不交且非空的开闭子集 U 和 V，使得 \(X - \{x\} = U \cup V\)。不妨设 \(S \subset V\)。有 \(\overline{U} = U \cup \{x\}\) 和 \(\overline{V} = V \cup \{x\}\)，且这些空间是连通的（III, p.~278, 命题 17, a)）。由本命题的断言 c)，存在 \(y \in U\)，使得 \(X - \{y\}\) 连通，这与包含关系 \(N \subset S \subset V\) 矛盾。

引理 7. --- 设 T 是一个局部连通的拓扑空间。设 \(\mathcal{U}\) 是 T 的开子集的有向递增族，其并为 T。对于 \(t \in T\) 和 \(U \in \mathcal{U}\)，记 \(U_{t}\) 为 t 在 U 中的连通分支，若 \(t \in U\)；置 \(U_{t} = \varnothing\)，若 \(t \notin U\)。对所有 \(t \in T\)，\(\bigcup_{U \in \mathcal{U}} U_{t}\) 是 t 在 T 中的连通分支。

对 \(t \in \mathrm{T}\)，置 \(\mathrm{C}_t = \bigcup_{\mathrm{U} \in \mathcal{U}} \mathrm{U}_t\)。有 \(t \in \mathrm{C}_t\)。集合 \(\mathrm{U}_t\) 是开且连通的，\(\mathrm{C}_t\) 也如此，因为有 \(t \in \mathrm{U}_t\)（当 \(\mathrm{U}_t \neq \varnothing\) 时）（TG, I, p.~81, 命题 2）。设 \(u\) 和 \(v\) 是 T 中满足 \(\mathrm{C}_u \cap \mathrm{C}_v \neq \varnothing\) 的点。令 \(t\) 为 \(\mathrm{C}_u \cap \mathrm{C}_v\) 中的一点；取 \(\mathrm{U} \in \mathcal{U}\) 使得 \(t \in \mathrm{U}_u\)，并取 \(\mathrm{V} \in \mathcal{U}\) 使得 \(v \in \mathrm{V}_v\)。由于 \(\mathcal{U}\) 按包含关系过滤递增，存在 \(\mathrm{W} \in \mathcal{U}\)，包含 \(\mathrm{U} \cup \mathrm{V}\)。于是 \(\mathrm{W}_u \cap \mathrm{W}_v \neq \varnothing\)，从而 \(\mathrm{W}_u = \mathrm{W}_v\)。更一般地，都有 \(\mathrm{W}_u' = \mathrm{W}_v'\)，对所有 \(\mathrm{W}' \in \mathcal{U}\) 满足 \(\mathrm{W} \subset \mathrm{W}'\) 的情形，从而 \(\mathrm{C}_u = \mathrm{C}_v\)。这表明，形如 \(\mathrm{C}_t\) 的集合构成 T 到连通开子集的一个划分。因此，对每个 \(t \in \mathrm{T}\)，\(\mathrm{C}_t\) 是 \(t\) 在 T 中的连通分支。

定理 2. --- 设 X 是一个具有可数稠密子集的紧连通拓扑空间。设 a、b 是 X 中彼此不同的点。下列条件等价：

\begin{enumerate}
\def\labelenumi{(\roman{enumi})}
\item
  存在一个同胚 \(f\colon\mathbf{I}\to\mathbf{X}\)，使得 \(f(0)=a\) 且 \(f(1)=b\)；
\item
  每个包含 \(\{a, b\}\) 的 X 的连通子集都等于 X；
\item
  空间 X 是局部连通的，并且每个包含 \(\{a,b\}\) 的 X 的连通紧子集都等于 X；
\item
  对每个 \(x \in X - \{a, b\}\)，空间 \(X - \{x\}\) 都不连通。
\end{enumerate}

断言 (i) 蕴含其余所有断言。

令 \(x\) 为 \(X - \{a, b\}\) 中的一点。由于 \(X - \{x\}\) 包含 \(\{a, b\}\)，断言 (ii) 蕴含它不是 \(X\) 的连通子集，故 (ii) 蕴含 (iv)。

我们来证明 (iii) 蕴含 (iv)。假设 (iii) 的假设成立。令 \(x \in X - \{a, b\}\)，并假设 \(X - \{x\}\) 是连通的。令 T 为 \(X - \{x\}\) 这个空间，令 \(\mathcal{U}\) 为由形如 \(X - V\) 的 T 的子集组成的集合，其中 V 是 \(x\) 的一个紧邻域。由引理 7，存在 \(x\) 的一个紧邻域 V，使得 \(a\) 和 \(b\) 属于 \(X - V\) 的同一个连通分支。特别地，它们属于 \(X - \mathring{V}\) 的同一个连通分支；这个分支是 X 的一个紧连通子集，并且不同于 X，这与假设 (iii) 矛盾。

还需证明断言 (iv) 蕴含条件 (i)。

令 \(x\) 为 \(X - \{a, b\}\) 中的一点，令 U、V 为 \(X - \{x\}\) 的彼此不交且非空的开闭子集，使得 \(X - \{x\} = U \cup V\)。由 III, p.~278, 命题 17, c)，U（分别为 V）中存在一点，其在 X 中的补集是连通的。由于 X 中只有两个这样的点，即 \(a\) 和 \(b\)，其中一个——不妨设为 \(a\)——包含于 U，另一个包含于 V。由上引文，U 的一个非空开闭子集 \(U'\) 包含 \(\{a, b\}\) 中的一点，从而包含 \(a\)。将此应用于 \(U - U'\)，可知 \(U = U'\)。因此 U 是连通的；它是 \(a\) 在 \(X - \{x\}\) 中的连通分支。同样，V 是 \(b\) 在 \(X - \{x\}\) 中的连通分支。

对每个 \(x \in X - \{a, b\}\)，令 \(U_x\) 和 \(V_x\) 分别表示 a 和 b 在 \(X - \{x\}\) 中的连通分支。根据上文，它们在 \(X - \{x\}\) 中是开闭的、彼此不交，且它们的并等于 \(X - \{x\}\)。

按如下方式定义 X 上的关系 \(\preccurlyeq\)：一方面，有 \(a \preccurlyeq x\) 和 \(x \preccurlyeq b\)，对所有 \(x \in \mathrm{X}\) 成立；另一方面，若 \(x\) 和 \(y\) 属于 \(\mathrm{X} - \{a, b\}\)，则定义 \(x \preccurlyeq y\) 当 \(x \in \overline{\mathrm{U}_y}\) 时。对于 \(x\) 和 \(y\)，若它们属于 \(\mathrm{X} - \{a, b\}\)，子集 \(\mathrm{V}_x\) 和 \(\mathrm{V}_y\) 有公共点 b，并且关系 \(x \preccurlyeq y\) 实际上等价于断言 \(\overline{\mathrm{U}_x} \subset \overline{\mathrm{U}_y}\)（III, p.~278, 命题 17, b)）。由此关系 \(\preccurlyeq\) 是 X 上的一个序关系。

设 \(x\) 和 \(y\) 是 X 中的点，且 \(x \preccurlyeq y\) 不成立。必有 \(x \neq a\) 且 \(y \neq b\)，并且 \(x \in V_y\)。若 \(x = b\) 或 \(y = a\)，则有 \(y \preccurlyeq x\)。现在假设 \(x\) 和 \(y\) 都异于 \(a\) 和 \(b\)。由于部分集 \(U_x\) 和 \(U_y\) 有公共点 a，有包含关系 \(\overline{V_x} \subset \overline{V_y}\)（上引文），从而取补集的闭包得到 \(\overline{U_y} \subset \overline{U_x}\)，进而 \(y \preccurlyeq x\)。因此，空间 X 中的关系 \(\preccurlyeq\) 是一个全序关系。此外，对每个 \(x \in X - \{a, b\}\)，有 \(U_x = ] \leftarrow, x[\) 和 \(V_x = ]x, \rightarrow[\)；对 \(x, y \in X - \{a, b\}\)，有 {]}x, y{[} = \(U_y \cap V_x\)。当 x 和 y 遍历 \(X - \{a, b\}\) 中的点时，集合 \(U_y \cap V_x\)、集合 \(U_y\) 以及集合 \(V_x\) 构成 X 上一个拓扑的基。记相应的拓扑空间为 \(\widetilde{X}\)，并令 i: \(X \to \widetilde{X}\) 为 X 的恒等映射。由于对每个 \(x \in X\)，集合 \(U_x\) 和 \(V_x\) 在 X 中是开的，映射 i 是连续的。

空间 \(\widetilde{\mathbf{X}}\) 是分离的。事实上，设 \(x\) 和 \(y\) 是 \(\widetilde{\mathbf{X}}\) 中满足 \(x \preccurlyeq y\) 的彼此不同的点。集合 \(]\leftarrow, y[\) 和 \(]x, \rightarrow[\) 是 \(x\) 和 \(y\) 的开邻域；若它们有公共点 \(z\)，则 \(]\leftarrow, z[\) 和 \(]z, \rightarrow[\) 是 \(x\) 和 \(y\) 的不交开邻域。由于 X 是紧的，映射 \(i\) 因而是同胚（TG, I, p.~63, 推论 2）。

因此，X 的一个可数稠密子集在 \(i\) 下的像与有序集 \((\mathrm{X}, \preccurlyeq)\) 的每个非空开区间相交。由此根据 III, p.~276 的命题 16，存在同构 \(c\)，它从有序集 \(\mathbf{I}\) 到 \((\mathrm{X}, \preccurlyeq)\)。根据该命题后的注，这个同构是从 \(\mathbf{I}\) 到拓扑空间 \(\widetilde{\mathrm{X}}\) 的同胚。于是映射 \(f = i^{-1} \circ c\) 是从 \(\mathbf{I}\) 到 X 的同胚，并将 0 映为 \(a\)、将 1 映为 \(b\)。

\subsection{单射道路}

命题 18. --- 设 X 是一个豪斯多夫拓扑空间。设 a 和 b 是 X 中属于同一道路连通分支的彼此不同的点。存在一条在 X 中连接 a 与 b 的单射道路。

设 \(f \colon \mathbf{I} \to \mathbf{X}\) 是连接 \(a\) 与 \(b\) 的 \(\mathbf{X}\) 中的一条道路。令 \(\mathcal{U}\) 为由开子集 \(\mathrm{U}\) 组成的集合，这些子集属于 {]}0,1{[}，并满足 \(f(x) = f(y)\) 对 U 的每个连通分支 {]}x, y{[} 成立。

引理 8. --- 按包含关系排序的集合 U 是归纳的。

设 V 是 U 的一个全序子集。我们来证明属于 V 的各集合的并 V 是 U 的一个元素，即，对于 V 的每个连通分支 {]}u, v{[}，都有 \(f(u) = f(v)\)。

令 \(x\) 为 \([u, v[\) 中的一点。令 \(\mathcal{V}_x\) 为这样的 \(U \in \mathcal{V}\) 组成的集合：\(x \in U\)；对于这样的 \(U\)，记 \([u_U, v_U[\) 为 \(x\) 在 \(U\) 中的连通分支。有 \(u_{U'} \leqslant u_U < v_U \leqslant v_{U'}\)，若 \(U\) 和 \(U'\) 是 \(\mathcal{V}_x\) 中满足 \(U \subset U'\) 的元素。由 III, p.~280 的引理 7，所有 \(U \in \mathcal{V}_x\) 的 \([u_U, v_U[\) 的并等于 {]}u, v{[}，故有 u = lim \(u_{U}\) 和 v = lim \(v_{U}\)，其中极限取遍有向集 \(V_{x}\)。由于映射 f 连续且拓扑空间 X 是豪斯多夫的，等式 \(f(u_{\mathrm{U}}) = f(v_{\mathrm{U}})\) 对所有 \(U \in V_{x}\) 蕴含 \(f(u) = f(v)\)，这正是所要证明的。

根据 E, III, p.~20, 定理 2，存在属于 \(\mathcal{U}\) 的一个开子集 U，它关于包含关系是极大的。

令 \(g \colon \mathbf{I} \to \mathrm{X}\) 按如下方式定义。若 \(t \notin \mathrm{U}\)，则置 \(g(t) = f(t)\)；否则，记 \(]u, v[\) 为 \(t\) 在 \(\mathrm{U}\) 中的连通分支，并置 \(g(t) = f(u) = f(v)\)，于是映射 \(g|[u, v]\) 是像为 \(f(u)\) 的常值映射。

命题 18 由下述引理推出。

引理 9. --- 存在连续的递增满射 \(u\colon\mathbf{I}\to\mathbf{I}\) 以及单射道路 \(c\colon\mathbf{I}\to\mathbf{X}\)，它连接 \(a\) 与 \(b\)，并使得 \(g=c\circ u\)。

我们先证明映射 \(g\) 连续。令 \(t\) 为 \(\mathbf{I}\) 中的一点。若 \(t \in \mathrm{U}\)，则 \(g\) 在 \(t\) 的一个邻域内为常值，因而在 \(t\) 处连续。假设 \(t \notin \mathrm{U}\)，令 \(\mathrm{W}\) 为 \(g(t)\) 在 \(\mathrm{X}\) 中的一个开邻域。令 \(\mathrm{V}\) 为 \(\mathbf{I}\) 中包含 \(t\) 且满足 \(f(\mathrm{V}) \subset \mathrm{W}\) 的一个开区间；这样的区间存在，因为 \(f\) 在 \(t\) 处连续。我们来证明 \(g(\mathrm{V}) \subset \mathrm{W}\)。令 \(x \in \mathrm{V}\)；若 \(x \notin \mathrm{U}\)，则有 \(g(x) = f(x)\)，因而 \(g(x) \in \mathrm{W}\)。于是设 \(x \in \mathrm{U}\)，并令 \(]u, v[\) 为 \(x\) 在 \(\mathrm{U}\) 中的连通分支。注意 \(]u, v[\) 不包含 \(t\)。因此，若 \(x < t\)，则有 \(x < v \leqslant t\)，从而 \(v \in \mathrm{V}\) 且 \(g(x) = g(v) = f(v) \in f(\mathrm{V}) \subset \mathrm{W}\)。同样可证 \(g(x) \in \mathrm{W}\)，若 \(x > t\)。于是 \(g\) 在 \(t\) 处连续。

设 \(u\) 和 \(v\) 是 \(\mathbf{I}\) 中满足 \(g(u) = g(v)\) 且 \(u < v\) 的点。我们来证明 \(]u, v[ \subset U\)。置 \(U' = U \cup ]u, v[\)；它是 \(]0, 1[\) 的一个开子集。

若 \(u\) 属于 U，记 \(u'\) 为 \(u\) 在 U 中的连通分支在 I 中的最大下界；置 \(u' = u\)，若 \(u\) 不属于 U。同样，若 \(v\) 属于 U，记 \(v'\) 为 \(v\) 在 U 中的连通分支在 I 中的最小上界；置 \(v' = v\)，若 \(v\) 不属于 U。于是 \(]u', v'[\) 是 U' 的一个连通分支，且有 \(f(u') = g(u) = g(v) = f(v')\)。由于 U' 中异于 \(]u', v'[\) 的连通分支都是 U 的连通分支，开集 U' 是 \(\mathcal{U}\) 的一个元素。由于 U 关于包含关系是 \(\mathcal{U}\) 的极大元素，有 U' = U，且 \(]u, v[\) 包含于 U。这证明了映射 \(g\) 在区间 \([u, v]\) 上为常值。因此，\(g\) 的纤维是 I 的区间；由于 g 连续，这些区间是闭的。

记 R 为与 g 相伴的等价关系，记 p 为从 I 到 I/R 的典范满射。存在唯一的连续映射 \(g'\) 从 I/R 到 X，使得 \(g = g' \circ p\)；该映射是单射。由 III, p.~276 的命题 16 的推论，存在递增、连续且满射的映射 u，使得 R 是与 u 相伴的等价关系。由于空间 I 是紧的而空间 X 是豪斯多夫的，映射 u 是闭映射，因而是严格的（I, p.~18, 例 2）。它通过取商定义了从 I/R 到 I 的同胚 \(u'\)。于是，从 I 到 X 的映射 \(g' \circ (u')^{-1}\) 是一条起点为 a、终点为 b 的单射道路。

\subsection{道路的提升}

定理 3. --- 设 I 是 R 的一个区间，X 是一个非空分离拓扑空间。设 p: X → I 是连续、开且逆紧的映射，其纤维全不连通（TG, I, p.~83）。映射 p 是满射。对 X 中的每个点 x，存在 p 的一个连续截面 s，使得 s(p(x)) = x。

集合 \(p(\mathrm{X})\) 是 I 的一个开、闭（TG, I, p.~72, 命题 1）且非空的子集。由于 I 是连通的，有 \(p(\mathrm{X}) = \mathrm{I}\)；因此映射 \(p\) 是满射。

对每个 \((a,b)\in\mathrm{I}\times\mathrm{I}\) 满足 \(a\leqslant b\) 的，记 \(F_{a,b}\) 为由 \((y,z)\in\overline{p}^{-1}(a)\times\overline{p}^{-1}(b)\) 组成的集合，其中 y 和 z 属于 \(\overline{p}^{-1}([a,b])\) 的同一个连通分支。

引理 10. --- 设 \(a\) 和 \(b\) 是 I 中满足 \(a \leqslant b\) 的点。

\begin{enumerate}
\def\labelenumi{\alph{enumi})}
\item
  \(F_{a,b}\) 在 \(\overline{p}^{-1}(a)\times\overline{p}^{-1}(b)\) 中是闭的。
\item
  有 \(\operatorname{pr}_1(\mathrm{F}_{a,b}) = \overline{p}^{-1}(a)\) 且 \(\operatorname{pr}_2(\mathrm{F}_{a,b}) = \overline{p}^{-1}(b)\)。
\item
  设 \(c \in I\) 满足 \(b \leqslant c\)。若 \((y,z)\) 属于 \(F_{a,b}\) 且 \((z,t)\) 属于 \(F_{b,c}\)，则 \((y,t)\) 属于 \(F_{a,c}\)。
\end{enumerate}

集合 \(\overline{p}^{-1}([a,b])\) 是紧的（TG, I, p.~77, 命题 7）。因此，\((y,z)\in\overline{p}^{-1}(a)\times\overline{p}^{-1}(b)\) 属于 \(F_{a,b}\) 的充要条件是：\(\overline{p}^{-1}([a,b])\) 中每个包含 \(y\) 的开闭子集也包含 \(z\)（TG, II, p.~32, 命题 6）。

置 \(\mathrm{Y} = \overline{p}^{-1}([a,b])\)。对于 Y 的每个既开又闭的子集 U，集合 \(((\mathrm{U} \times \mathrm{U}) \cup ((\mathrm{Y} - \mathrm{U}) \times (\mathrm{Y} - \mathrm{U}))) \cap (\overline{p}^{-1}(a) \times \overline{p}^{-1}(b))\) 在 \(\overline{p}^{-1}(a) \times \overline{p}^{-1}(b)\) 中是闭的。这些集合的交等于 \(\mathrm{F}_{a,b}\)，故得出 \(a\)。

令 \(y \in \overline{p}^{-1}(a)\)。记 \(\mathcal{U}\) 为 Y 中 y 的开闭邻域组成的集合。由 p 取子空间诱导的从 Y 到 {[}a, b{]} 的映射是开且逆紧的（I, p.~17, 命题 8），因而也是闭的。于是，对每个 \(U\) 属于 \(\mathcal{U}\) 的，\(p(U)\) 是 {[}a, b{]} 的非空开闭子集；由于区间 {[}a, b{]} 是连通的，有 \(p(U) = [a, b]\)，特别地 \(U \cap \overline{p}^{-1}(b) \neq \varnothing\)。因此，\((U \cap \overline{p}^{-1}(b))_{U \in \mathcal{U}}\) 是紧空间 \(\overline{p}^{-1}(b)\) 的一个由非空闭子集组成的递减过滤族。该族的交非空（TG, I, p.~59）；令 z 为其中一个元素。有 \((y, z) \in \mathrm{F}_{a,b}\)。我们证明了关系 \(\mathrm{pr}_{1}(\mathrm{F}_{a,b}) = \overline{p}^{-1}(a)\)。将 p、I、a、b 分别替换为 -p、-I、-b、-a，即可由此得到关系 \(\mathrm{pr}_{2}(\mathrm{F}_{a,b}) = \overline{p}^{-1}(b)\)。

在 \(c\) 的假设下，\((y,t)\) 属于 \(\overline{p}^{-1}(a)\times\overline{p}^{-1}(c)\)。集合 \(\{y,z\}\) 包含于 \(\overline{p}^{-1}([a,b])\) 的一个连通子集 C 中，集合 \(\{z,t\}\) 包含于 \(\overline{p}^{-1}([b,c])\) 的一个连通子集 C' 中。于是，C \(\cup\) C' 是 \(\overline{p}^{-1}([a,c])\) 的连通子集（TG, I, p.~81, 命题 2），并包含 \(\{y,t\}\)，从而有关系 \((y,t)\in\mathrm{F}_{a,c}\)。

回到定理 3 的证明。映射 p 的每个纤维都是紧的（TG, I, p.~77, 命题 7）。由蒂霍诺夫定理（TG, I, p.~63, 定理 3），积空间 \(K = \prod_{a \in I} \overline{p}^{-1}(a)\) 是紧的。令 \(K'\) 为 K 中的元素 \((y_{a})_{a \in I}\) 组成的集合：它们满足 \(y_{p(x)}\) 等于 x，并且都有 \((y_{a}, y_{b}) \in F_{a,b}\)，对 I 中的每一对 \((a, b)\) 当 a \textless{} b 时成立。定理 3 由下述引理推出。

引理 11. --- a) 集合 K' 非空。

\begin{enumerate}
\def\labelenumi{\alph{enumi})}
\setcounter{enumi}{1}
\tightlist
\item
  设 \((s_{a})_{a\in\mathrm{I}}\) 是 \(K'\) 的一个元素。映射 \(s\colon a\mapsto s_{a}\) 从 I 到 X，是 p 的一个连续截面，且满足 \(s(p(x))=x\)。
\end{enumerate}

对于包含点 \(p(x)\) 的 I 的每个有限子集 S，记 \(K_{S}\) 为 K 中的元素 \((y_{a})_{a\in\mathrm{I}}\) 组成的集合，它们满足关系 \(y_{p(x)} = x\)，并且都有 \((y_{a}, y_{b}) \in \mathrm{F}_{a,b}\)，对 S 中的每一对 \((a, b)\) 当 a \textless{} b 时成立。集合 \(K_{S}\) 在 K 中是闭的（引理 10, a)），并构成 K 的一个递减过滤族，其交为 \(K'\)。

要证明 \(K'\) 非空，只需证明：对于包含点 \(p(x)\) 的 I 的每个有限子集 S，集合 \(K_{S}\) 非空（TG, I, p.~59）。

设 S 是这样的一个子集；将其元素排列成严格递增序列 \((a_{1},\ldots,a_{n})\)，并令 i 为满足 \(p(x)=a_{i}\) 的整数。置 \(y_{a_{i}}=x\)。由引理 10, b)，可以通过归纳构造元素 \(y_{a_{j}}\in\overline{p}^{-1}(a_{j})\)，其中 \(i<j\leqslant n\)，并通过逆向归纳构造元素 \(y_{a_{j}}\in\overline{p}^{-1}(a_{j})\)，其中 \(1\leqslant j<i\)，使得有 \((y_{a_{j}},y_{a_{j+1}})\in\mathrm{F}_{a_{j},a_{j+1}}\)，对每个满足 \(1\leqslant j<n\) 的整数 j 成立。由引理 10, c)，有 \((y_{a},y_{b})\in\mathrm{F}_{a,b}\)，对 S 中的每一对 \((a,b)\) 满足 a\textless b。由于映射 p 是满射，可以对每个 \(a\in I-S\) 选取一个元素 \(y_{a}\in\overline{p}^{-1}(a)\)。如此构造的族 \((y_{a})_{a\in I}\) 属于 \(K_{S}\)，因此 \(K_{S}\) 非空。

我们来证明 \(b\)。根据 \(K'\) 的定义，\(s\) 是 \(p\) 的一个截面，且满足 \(s(p(x)) = x\)。令 \(a \in I\)。我们来证明 \(s\) 在点 \(a\) 处连续。令 U 为 X 中 \(s_a\) 的一个开邻域。由于 \(\overline{p}^{-1}(a)\) 是紧空间（TG, I, p.~77, 命题 7）且全不连通，\(s_a\) 在 \(\overline{p}^{-1}(a)\) 中有一个包含于 U 的开闭邻域 C（TG, II, p.~32, 推论）。集合 C 和 \(\overline{p}^{-1}(a) - C\) 在 \(\overline{p}^{-1}(a)\) 中是闭的，因而是紧的，并且彼此不交。由于 X 是分离的，它们有不交的开邻域 V 和 \(V'\)。集合 \((V \cap U) \cup V'\) 是 X 中纤维 \(\overline{p}^{-1}(a)\) 的一个开邻域；由于映射 \(p\) 是闭的（TG, I, p.~72, 命题 1），\((V \cap U) \cup V'\) 包含形如 \(\overline{p}^{-1}(J)\) 的集合，其中 \(J \subset I\) 是包含 \(a\) 的开区间（I, p.~75, 引理）。置 \(W = V \cap U \cap \overline{p}^{-1}(J)\)。集合 W 在 \(\overline{p}^{-1}(J)\) 中是开的；在 \(\overline{p}^{-1}(J)\) 中也为闭的，因为有 \(\overline{p}^{-1}(J) - W = V' \cap \overline{p}^{-1}(J)\)。令 \(b \in J\)。以 \(a\) 和 \(b\) 为端点的 I 的闭区间包含于 J 中。根据假设，\((s_a, s_b)\) 属于 \(F_{a,b}\)，当 \(a \leqslant b\) 时；\((s_b, s_a)\) 属于 \(F_{b,a}\)，当 \(b \leqslant a\) 时。因此，\(\overline{p}^{-1}(J)\) 中存在包含 \(\{s_a, s_b\}\) 的连通子集。于是，\(s_b\) 属于 \(\overline{p}^{-1}(J)\) 中每个包含 \(s_a\) 的开闭子集，特别地属于 W；从而 \(s_b\) 属于 U。因此有 \(s(J) \subset U\)，这证明了 \(s\) 在点 \(a\) 处连续。

推论. --- 设 X 和 B 是拓扑空间，\(p: X \to B\) 是连续、开、逆紧且分离的映射，其纤维全不连通。设 I 是 R 的一个区间，\(f: I \to B\) 是连续映射，a 是 I 中的一点，x 是 X 中满足 \(f(a) = p(x)\) 的一点。存在连续映射 \(g: I \to X\)，使得 \(p \circ g = f\) 且 \(g(a) = x\)。

置 \(X' = I \times_{B} X\)，并记 \(p'\colon X'\to I\) 和 \(f'\colon X'\to X\) 为典范投影。映射 \(p'\) 是连续、开、逆紧且分离的（I, p.~17, 命题 8 及 p.~27, 命题 4）。由于空间 I 是分离的，空间 \(X'\) 也是分离的（I, p.~26, 注 3）。\(p'\) 的纤维全不连通（I, p.~10, 推论, a)）。由定理 3，存在连续截面 \(s'\)，它是 \(p'\) 的截面，并取值 \((a,x)\) 于 \(a\) 处。从 I 到 X 的映射 \(g = f'\circ s'\) 是连续的；有 \(p\circ g = p\circ f'\circ s' = f\circ p'\circ s' = f\) 且 \(g(a) = x\)，故得推论。

定理 4. --- 设 X 是一个拓扑空间，G 是一个离散群，在 X 上作逆紧作用，p: X → X/G 是典范映射。设 I 是 R 的一个区间，f: I → X/G 是连续映射，a 是 I 中的一点，x 是 X 中满足 \(f(a) = p(x)\) 的一点。存在连续映射 \(\varphi\)：I → X，使得 \(p \circ \varphi = f\) 且 \(\varphi(a) = x\)。

我们先处理 I 是 R 的一个有界闭区间的情形。

根据 TG, III, p.~29, 命题 3，空间 X 是分离的。

令 \(y\) 为 X/G 中的一点，并令 \(x \in X\) 满足 \(y = p(x)\)。\(K_x\) 是 \(x\) 的稳定子，是 G 的有限子群；此外，存在 X 中的邻域 \(U_x\)，其为点 \(x\) 的邻域，以及 X/G 中的邻域 \(V_y\)，其为点 \(y\) 的邻域，使得 \(U_x\) 在 \(K_x\) 下稳定，\(gU_x \cap U_x = \varnothing\)，若 \(g \notin K_x\)，且 \(p\) 诱导从 \(U_x/K_x\) 到 \(V_y\) 的同胚。另外，对每个 \(g \in G\)，\(p\) 诱导从 \(gU_x\) 到 \(V_y\) 的同胚。由于 I 是紧的，存在整数 \(m\) 和 \(n\)，满足 \(m \leqslant 0 \leqslant n\)，以及 I 中元素的有限序列 \((a_i)_{m \leqslant i \leqslant n}\)，使得 \(a_0 = a\)、\(I = {[}a_m,a_n{]}\)，并且对每个 \(i \in \{m, \ldots, n-1\}\)，\(f({[}a_i,a_{i+1}{]})\) 包含于按上述方式构造的 X/G 的一个开集 \(V_{y_i}\) 中。

令 \(x_0\) 为 \(\overline{p}^{-1}(y_0)\) 中满足 \(x \in \mathrm{U}_{x_0}\) 的唯一元素。令 \(q_0\) 为通过取商从 \(\mathrm{U}_{x_0}\) 到 \(\mathrm{U}_{x_0}/\mathrm{K}_{x_0}\) 的典范映射。映射 \(p\) 诱导同胚 \(i_0\)，从 \(\mathrm{U}_{x_0}/\mathrm{K}_{x_0}\) 到 \(\mathrm{V}_{y_0}\)，使得 \(i_0 \circ q_0 = p \mid \mathrm{U}_{x_0}\)。映射 \(q_0\) 是逆紧的（TG, III, p.~29, 命题 3）、开的（TG, I, p.~31, 例 1）且分离的，因为 X 是分离的。它的纤维全不连通，因为它们是有限集。由 III, p.~286 的推论，存在连续映射 \(\varphi_0\colon [a_0,a_1]\to \mathrm{U}_{x_0}\)，使得 \(p\circ \varphi_0 = f\mid [a_0,a_1]\)。

类似地，通过对整数 \(i \in \{0, \ldots, n-1\}\) 归纳，我们构造 \(x_i \in \overline{p}^{-1}(y_i)\) 以及连续映射 \(\varphi_i: [a_i, a_{i+1}] \to X\)，其像包含于 \(\mathrm{U}_{x_i}\) 中，并满足 \(p \circ \varphi_i = f \mid [a_i, a_{i+1}]\)，且满足 \(\varphi_i(a_{i+1}) = \varphi_{i+1}(a_{i+1})\)，当 \(0 \leqslant i < n-1\) 时。

类似地，通过对整数 \(i \in \{m, \ldots, -1\}\) 作逆向归纳，构造 \(x_i \in \overline{p}^{-1}(y_i)\) 以及连续映射 \(\varphi_i: [a_i, a_{i+1}] \to X\)，其像包含于 \(U_{x_i}\) 中，并满足 \(p \circ \varphi_i = f \mid [a_i, a_{i+1}]\)，且满足 \(\varphi_i(a_{i+1}) = \varphi_{i+1}(a_{i+1})\)，当 \(m \leqslant i < 0\) 时。

存在唯一映射 \(\varphi\colon\mathrm{I}\to\mathrm{X}\)，它与 \(\varphi_{i}\) 在 \([a_{i},a_{i+1}]\) 上一致，对满足 \(m\leqslant i<n\) 的 i 成立。它是连续的（TG, I, p.~19, 命题 4），是 \(f\) 到 X 的连续提升，并满足 \(\varphi(a)=x\)。

当 I 紧时，定理得证。在一般情形下，存在序列 \((a_{n})_{n\in\mathbf{N}}\) 和 \((b_{n})_{n\in\mathbf{N}}\)，使得 \((a_{n})\) 以 \(\inf(I)\) 为极限且最终为常值，\((b_{n})\) 以 \(\sup(I)\) 为极限且最终为常值，\(a=a_{0}=b_{0}\)，并且当 I 有最小元（分别有最大元）时，\((a_{n})\)（分别 \((b_{n})\)）为常值。根据上文，对每个整数 \(n\in N\)，存在连续提升 \(\varphi_{n}\) 为 \(f\mid[a_{n},a_{n+1}]\) 到 X 的提升，以及连续提升 \(\varphi_{n}^{\prime}\) 为 \(f\mid[b_{n+1},b_n]\) 到 X 的提升，使得 \(\varphi_{0}(a_{0})=\varphi_{0}^{\prime}(b_{0})=x\)，且 \(\varphi_{n+1}(a_{n+1})=\varphi_{n}(a_{n+1})\)、\(\varphi_{n+1}^{\prime}(b_{n+1})=\varphi_{n}^{\prime}(b_{n+1})\)。存在唯一映射 \(\varphi\colon I\to X\)，它与\(\varphi_{n}\) 在 \([a_{n},a_{n+1}]\) 上一致，并与\(\varphi_{n}^{\prime}\) 在 \([b_{n+1},b_{n}]\) 上一致，对每个 \(n\in N\) 均如此。它是连续的（上引文），是 f 到 X 的连续提升，并满足 \(\varphi(a)=x\)。定理得证。

例. --- 1) 设 X 是一个分离拓扑空间，G 是赋予离散拓扑的有限群，且在 X 上连续作用。于是该作用是逆紧的（TG, III, p.~28, 命题 2）。定理 4 的断言于是直接由定理 3 的推论得到。

\begin{enumerate}
\def\labelenumi{\arabic{enumi})}
\setcounter{enumi}{1}
\tightlist
\item
 设 n 是一个满足 \(\geqslant 0\) 的整数。记 \(P_{n}\) 为次数为 n 的首一多项式 \(P \in C[X]\) 组成的集合，并在其上赋予这样的拓扑：映射 \((c_{0}, \ldots, c_{n-1}) \mapsto X^{n} + c_{n-1}X^{n-1} + \cdots + c_{0}\) 是从 \(C^{n}\) 到 \(P_{n}\) 的同胚。由从 \(C^{n}\) 到 \(P_{n}\) 的映射 p，其中 \(p(z_{1}, \ldots, z_{n}) = (X - z_{1}) \ldots (X - z_{n})\)，是连续的。对称群 \(S_{n}\) 通过置换因子作用在 \(C^{n}\) 上，而 p 通过取商定义了从 \(C^{n}/S_{n}\) 到 \(P_{n}\) 的同胚（TG, VIII, p.~22, 命题 1，I, p.~23, 推论 1 及 TG, VIII, p.~20）。因此得到下述论断：
\end{enumerate}

设 I 是 R 的一个区间，\((c_{0},\ldots,c_{n-1})\) 是从 I 到 C 的连续映射序列，a 是 I 中的一点，\((z_{1},\ldots,z_{n})\) 是一个复数序列，满足 \((\mathrm{X}-z_{1})\ldots(\mathrm{X}-z_{n})=\mathrm{X}^{n}+c_{n-1}(a)\mathrm{X}^{n-1}+\cdots+c_{0}(a)\)。存在一个从 I 到 C 的连续映射序列 \((\lambda_{1},\ldots,\lambda_{n})\)，使得 \(\lambda_{i}(a)=z_{i}\) 对 \(1\leqslant i\leqslant n\) 成立，并且有 \((\mathrm{X}-\lambda_{1}(t))\ldots(\mathrm{X}-\lambda_{n}(t))=\mathrm{X}^{n}+c_{n-1}(t)\mathrm{X}^{n-1}+\cdots+c_{0}(t)\) 对所有 \(t\in I\) 成立。

\section{庞加莱群胚}

\subsection{庞加莱群胚}

定义 1. --- 设 X 是一个拓扑空间，\(c_{0}\) 和 \(c_{1}\) 是 X 中的道路，\(\sigma\colon\mathbf{I}\times\mathbf{I}\to\mathbf{X}\) 是连接 \(c_{0}\) 与 \(c_{1}\) 的同伦。称 \(\sigma\) 为严格同伦，若映射 \(s\mapsto\sigma(0,s)\) 和 \(s\mapsto\sigma(1,s)\) 为常值。

如果存在连接 \(c_{0}\) 与 \(c_{1}\) 的严格同伦，就称 X 中的道路 \(c_{0}\) 和 \(c_{1}\) 严格同伦。

两条严格同伦的道路具有相同的起点和终点。

例. --- 设 c 是 X 中的一条道路，\(\varphi\colon I\to I\) 是连续映射，满足 \(\varphi(0)=0\) 且 \(\varphi(1)=1\)。道路 c 与 \(c\circ\varphi\) 严格同伦。事实上，由 \(\sigma\colon I\times I\to X\) 和 \(\sigma(t,s)=c((1-s)t+s\varphi(t))\) 定义的映射是连接 c 与 \(c\circ\varphi\) 的严格同伦。

设 X 是一个拓扑空间。回忆（参见 III, p.~257），\(\Lambda(X)\) 表示道路空间 \(\mathcal{C}_c(\mathbf{I};\mathrm{X})\)，并且对于 \(x\)、\(y \in X\)，\(\Lambda_{x,y}(X)\) 是 \(\Lambda(X)\) 的子空间，由起点为 x、终点为 y 的道路组成。对于 \(x\)、\(y\)，集合 \(\Lambda_{x,y}(X)\) 是 X 的道路空间的一个划分。根据（III, p.~257, 注 2）从 \(x\)、\(y \in X\)、\(\mathcal{C}(\mathbf{I} \times \mathbf{I};\mathrm{X})\) 到 \(\mathcal{C}(\mathbf{I};\Lambda(\mathrm{X}))\) 的典范双射，严格同伦对应于这样的道路 \(c: \mathbf{I} \to \Lambda(X)\)，其像包含于形如 \(\Lambda_{x,y}(\mathrm{X})\) 的子空间中。因此，关系“道路 \(c_0\) 和 \(c_1\) 严格同伦”是 \(\Lambda(\mathrm{X})\) 上的等价关系（III, p.~259, 命题 5），而该关系的等价类是 X 的道路空间的子空间 \(\Lambda_{x,y}(\mathrm{X})\) 的道路连通分支。记 \(\varpi_{x,y}(\mathrm{X})\) 为集合 \(\pi_0(\Lambda_{x,y}(\mathrm{X}))\)，并称连接 \(x\) 与 y 的道路类为 \(y\)、\(\varpi_{x,y}(\mathrm{X})\) 的任一元素。

定义 2. --- 我们称 X 的圈空间为 \(\Omega(X)\)，它是 \(\Lambda(X)\) 的由 X 中的圈组成的子空间（III, p.~256, 定义 1）。

记 \(\Omega_x(\mathrm{X})\) 为集合 \(\Lambda_{x,x}(\mathrm{X})\)。\(\Omega_x(\mathrm{X})\) 的元素称为 X 中 \(x\) 处的圈，\(\varpi_{x,x}(\mathrm{X})\) 的元素称为 X 中 \(x\) 处的圈类。映射 \(e_x: \mathbf{I} \to \mathbf{X}\) 以 \(x\) 为像中的唯一点，是一个圈，称为 \(x\) 处的常值圈；其严格同伦类记为 \(\varepsilon_x\)。映射 \(x \mapsto e_x\) 从 X 到 \(\Lambda(\mathrm{X})\) 是连续的（III, p.~257, 命题 1）。

设 X 是一个拓扑空间，x、y、z 是 X 中的点。由连续映射 \(c \mapsto \overline{c}\)（从 \(\Lambda_{x,y}(X)\) 到 \(\Lambda_{y,x}(X)\)）（III, p.~258, 推论）取道路连通分支，得到从 \(\varpi_{x,y}(X)\) 到 \(\varpi_{y,x}(X)\) 的映射，记为 \(\gamma \mapsto \overline{\gamma}\)。若 \(\gamma \in \varpi_{x,y}(X)\)，则称 \(\overline{\gamma}\) 为道路类 \(\gamma\) 的逆。

同样，在将集合 \(\pi_0(\Lambda_{x,y}(\mathrm{X})) \times \pi_0(\Lambda_{y,z}(\mathrm{X}))\) 与 \(\pi_0(\Lambda_{x,y}(\mathrm{X}) \times \Lambda_{y,z}(\mathrm{X}))\) 识别之后（III, p.~260, 命题 6），由连续映射 \((c,d) \mapsto c * d\) 从 \(\Lambda_{x,y}(\mathrm{X}) \times \Lambda_{y,z}(\mathrm{X})\) 到 \(\Lambda_{x,z}(\mathrm{X})\)（III, p.~258, 推论）取道路连通分支后得到映射 \(C_{x,y,z} \colon \varpi_{x,y}(\mathrm{X}) \times \varpi_{y,z}(\mathrm{X}) \to \varpi_{x,z}(\mathrm{X})\)。对于 \(\gamma \in \varpi_{x,y}(\mathrm{X})\) 和 \(\delta \in \varpi_{y,z}(\mathrm{X})\)，记 \(\gamma \delta\) 为严格同伦类 \(C_{x,y,z}(\gamma,\delta)\)。称其为可拼接道路类 \(\gamma\) 与 \(\delta\) 的复合。

有 \(\overline{\overline{\gamma}}=\gamma\) 以及 \(\overline{\gamma\delta}=\overline{\delta}\overline{\gamma}\)。

命题 1. --- 设 X 是一个拓扑空间，x、y、z、u 是 X 中的点，\(\gamma_{1} \in \varpi_{x,y}(X)\)、\(\gamma_{2} \in \varpi_{y,z}(X)\)、\(\gamma_{3} \in \varpi_{z,u}(X)\) 是道路类。有

\begin{enumerate}
\def\labelenumi{(\arabic{enumi})}
\tightlist
\item
\end{enumerate}

\[
\varepsilon_ {x} \gamma_ {1} = \gamma_ {1} \varepsilon_ {y} = \gamma_ {1},\tag{2}
\]

\[
\gamma_ {1} \overline {{\gamma}} _ {1} = \varepsilon_ {x}, \qquad \overline {{\gamma}} _ {1} \gamma_ {1} = \varepsilon_ {y},\tag{3}
\]

\[
(\gamma_ {1} \gamma_ {2}) \gamma_ {3} = \gamma_ {1} (\gamma_ {2} \gamma_ {3}).
\]

令 \(c_{1} \in \Lambda_{x,y}(\mathrm{X})\)、\(c_{2} \in \Lambda_{y,z}(\mathrm{X})\)、\(c_{3} \in \Lambda_{z,u}(\mathrm{X})\) 分别为 \(\gamma_{1}\)、\(\gamma_{2}\) 和 \(\gamma_{3}\) 的代表元。令 \(\varphi: \mathbf{I} \to \mathbf{I}\) 为下式定义的函数：

\[
\varphi (t) = \left\{ \begin{array}{l l} t / 2 & \text { for } 0 \leqslant t \leqslant 1 / 2, \\ t - 1 / 4 & \text { for } 1 / 2 \leqslant t \leqslant 3 / 4, \\ 2 t - 1 & \text { for } 3 / 4 \leqslant t \leqslant 1. \end{array} \right.\tag{4}
\]

函数 \(\varphi\) 在三个区间 \([0,1/2]\)、\([1/2,3/4]\) 和 \([3/4,1]\) 上分别是仿射的，因而连续。有 \(\varphi(0) = 0\) 和 \(\varphi(1) = 1\)。由 III, p.~256 定义道路拼接的公式 (1)，可得

\[
c _ {1} * (c _ {2} * c _ {3}) = ((c _ {1} * c _ {2}) * c _ {3}) \circ \varphi .
\]

根据 III, p.~289 的例，路径 \(c_{1} * (c_{2} * c_{3})\) 与 \((c_{1} * c_{2}) * c_{3}\) 严格同伦，故等式 (3) 成立。

同样，由下式定义的函数 \(\psi\colon\mathbf{I}\to\mathbf{I}\)

\[
\psi (t) = \left\{ \begin{array}{l l} 2 t & \text { for } 0 \leqslant t \leqslant 1 / 2, \\ 1 & \text { for } 1 / 2 \leqslant t \leqslant 1 \end{array} \right.\tag{5}
\]

是连续的，并满足 \(\psi(0)=0,\psi(1)=1\)。等式 \(\gamma_{1}\varepsilon_{y}=\gamma_{1}\) 由下述事实得到：

\[
c _ {1} * e _ {y} = c _ {1} \circ \psi
\]

而等式 \(\varepsilon_{x}\gamma_{1}=\gamma_{1}\) 的证明相同，故得 (1)。

由下式定义的映射 \(\sigma\colon\mathbf{I}\times\mathbf{I}\to\mathrm{X}\)

\[
\sigma (t, s) = \left\{ \begin{array}{l l} c _ {1} (2 t s) & \text { for } 0 \leqslant t \leqslant 1 / 2, \\ c _ {1} (2 (1 - t) s) & \text { for } 1 / 2 \leqslant t \leqslant 1 \end{array} \right.\tag{6}
\]

是连续的；它是连接道路 \(e_{x}\) 与道路 \(c_{1}*\overline{c_{1}}\) 的严格同伦，故得到 (2) 中的第一个等式。第二个等式由第一个等式以及对每条道路 c 都有 \(c=\overline{\overline{c}}\) 这一事实得到。

注 1. --- 设 X 是一个拓扑空间。设 n 是满足 \(\geqslant 1\) 的整数，\((c_{1},\ldots,c_{n})\) 是 X 中的道路序列，对 \(c_{i}\)，\(c_{i+1}\) 与 \(1 \leqslant i \leqslant n-1\) 可拼接（这样的序列称为可拼接道路序列）。记道路 c 为

\[
c _ {1} * \left(c _ {2} * \left(\dots * \left(c _ {n - 1} * c _ {n}\right) \ldots\right)\right)
\]

以及由 \(c'\)（其中 \(c'(t) = c_i(nt - i + 1)\) 且 \(1 \leqslant i \leqslant n\)）定义的道路 \(t \in [\frac{i-1}{n}, \frac{i}{n}]\)。道路 \(c\) 与 \(c'\) 具有相同的像并且严格同伦：其中一条是另一条与一个固定 0 和 1 的 I 的同胚的复合（参见 III, p. 289, 例）。我们有时将道路 \(c_{1} * c_{2} * \cdots * c_{n}\) 记为 \(c'\)。

存在唯一的有向图 \(\varpi(\mathrm{X})\)，其顶点集为 X，连接点 \(x\) 与点 \(y\) 的箭集为 \(\varpi_{x,y}(\mathrm{X})\)，且映射 \(C_{x,y,z}\) 定义一个合成法则。由命题 1，\(\varpi(\mathrm{X})\) 是一个群胚（II, p.~162, 定义 4）。对每个 \(x \in \mathrm{X}\)，该群胚在顶点 x 处的单位元是像为 x 的常值圈的类。箭 \(x\)、\(x\)、\(\gamma\) 的逆是箭 \(\overline{\gamma}\)，我们也记为 \(\gamma^{-1}\)。特别地，合成法则 \(C_{x,x,x}\) 赋予每个 \(x \in \mathrm{X}\) 的集合 \(\varpi_{x,x}(\mathrm{X})\) 一个群结构；将此群记为 \(\pi_1(\mathrm{X}, x)\)。

定义 3. --- 设 X 是一个拓扑空间。群胚 \(\varpi(X)\) 称为 X 的庞加莱群胚或基本群胚。设 x 是 X 中的一点；x 处圈类组成的群 \(\pi_1(X,x)\) 称为 X 在点 x 处的庞加莱群或基本群。

设 \(\mathcal{U}\) 是 X 的一个子集族，其内部覆盖 X。X 中像包含于 \(\mathcal{U}\) 的某个成员中的道路类生成群胚 \(\varpi(\mathrm{X})\)（III, p.~272 的引理 4）。

群胚 \(\varpi(X)\) 的轨道（II, p.~162）与空间 X 的道路连通分支一致。特别地（上引文），有：

命题 2. --- 设 X 是一个拓扑空间。

\begin{enumerate}
\def\labelenumi{\alph{enumi})}
\item
  如果空间 X 是道路连通的，则对 \(\pi_1(\mathrm{X},x)\) 中的所有点 \(\pi_1(\mathrm{X},y)\) 和 \(x\)，群 \(y\) 与 \(\mathrm{X}\) 同构。
\item
  设 x 是 X 中的一点；下列条件等价：
\end{enumerate}

\begin{enumerate}
\def\labelenumi{(\roman{enumi})}
\item
  群 \(\pi_1(X,x)\) 是平凡的；
\item
  X 中起点为 x 且具有相同终点的两条道路严格同伦；
\item
  X 中每个起点为 x 的圈都与像为 x 的常值圈严格同伦。
\end{enumerate}

更确切地说，设 X 是一个道路连通的拓扑空间，x 和 y 是 X 中的点。对 \(\delta\) 的每个元素 \(\varpi_{x,y}(X)\)，映射 \(u_{\delta} \colon \gamma \mapsto \delta \gamma \delta^{-1}\) 从 \(\pi_1(\mathrm{X},y)\) 到 \(\pi_1(\mathrm{X},x)\) 是群同构，其逆同构为 \(u_{\delta^{-1}}\)。对于 \(\delta \in \varpi_{x,y}(\mathrm{X})\) 和 \(\gamma \in \pi_1(\mathrm{X},y)\)，置 \(^{\delta}\gamma = u_{\delta}(\gamma)\)。当 \(x = y\) 时，有 \(^{\delta}\gamma = \operatorname{Int}(\delta)(\gamma)\)，即 \(u_{\delta} = \operatorname{Int}(\delta)\)。

设 \(x, y, z\) 是 X 中的点。对于 \(\delta \in \varpi_{x,y}(X)\) 和 \(\eta \in \varpi_{y,z}(X)\)，有 \(u_{\delta \eta} = u_{\delta} \circ u_{\eta}\)，这也可对 \(\delta \eta \gamma = {}^{\delta}(\eta \gamma)\) 写成 \(\gamma \in \pi_1(X, z)\)。

设 \(x, y\) 是 X 中的点，\(\delta\)、\(\delta'\) 是 \(\varpi_{x,y}(\mathrm{X})\) 的元素。有

\[
u _ {\delta^ {\prime}} = u _ {\delta} \circ \operatorname{Int} (\delta^ {- 1} \delta^ {\prime}) = \operatorname{Int} (\delta^ {\prime} \delta^ {- 1}) \circ u _ {\delta}.\tag{7}
\]

注. --- 2) 设 X 是一个拓扑空间，C 是 X 的一个道路连通分支。若 x 和 y 是 C 中的点，则拓扑空间 \(\Lambda_{x,y}(\mathrm{C})\) 可与拓扑空间 \(\Lambda_{x,y}(\mathrm{X})\) 识别，从而集合 \(\varpi_{x,y}(\mathrm{C})\) 可与集合 \(\varpi_{x,y}(\mathrm{X})\) 识别。因此，基本群胚 \(\varpi(\mathrm{C})\) 可与 \(\varpi(\mathrm{X})\) 的、以 C 为点集的全子群胚识别。特别地，对 C 中的每个点 x，群 \(\pi_1(\mathrm{C}, x)\) 可与群 \(\pi_1(\mathrm{X}, x)\) 识别。

\begin{enumerate}
\def\labelenumi{\arabic{enumi})}
\setcounter{enumi}{2}
\tightlist
\item
  设 X 是一个拓扑空间，\(x\)、\(y\)、\(z\) 是 X 中的点。从 \((c,d)\mapsto c*d\) 到 \(\Lambda_{x,y}(\mathrm{X})\times \Lambda_{y,z}(\mathrm{X})\) 的映射 \(\Lambda_{x,z}(\mathrm{X})\) 是连续的（III, p.~258, 命题 2 的推论）。注意，当赋予集合 \(\varpi_{x,y}(\mathrm{X})\times \varpi_{y,z}(\mathrm{X})\to \varpi_{x,z}(\mathrm{X})\)、\(\varpi_{x,y}(\mathrm{X})\) 和 \(\varpi_{y,z}(\mathrm{X})\) 商拓扑时，由它诱导的合成映射 \(\varpi_{x,z}(\mathrm{X})\) 不一定连续（参见 TG, I, p.~35 及 III, p.~259）。然而，对每个 \(\gamma_0\in \varpi_{x,y}(\mathrm{X})\) 和每个 \(\delta_0\in \varpi_{y,z}(\mathrm{X})\)，部分映射 \(\gamma \mapsto \gamma \delta_0\) 从 \(\varpi_{x,y}(\mathrm{X})\) 到 \(\varpi_{x,z}(\mathrm{X})\)，以及 \(\delta \mapsto \gamma_0\delta\) 从 \(\varpi_{y,z}(\mathrm{X})\) 到 \(\varpi_{x,z}(\mathrm{X})\)，都是同胚。事实上，令 \(c_0\in \Lambda_{x,y}(\mathrm{X})\) 为类为 \(\gamma_0\) 的道路。映射 \(d\mapsto c_0*d\) 是从 \(\Lambda_{y,z}(\mathrm{X})\) 到 \(\Lambda_{x,z}(\mathrm{X})\) 的连续映射。映射 \(\delta \mapsto \gamma_0\delta\) 通过取商得到，因而连续。从 \(\delta^{\prime}\mapsto \gamma_{0}^{-1}\delta^{\prime}\) 到 \(\varpi_{x,z}(\mathrm{X})\) 的映射 \(\varpi_{y,z}(\mathrm{X})\) 同样如此。由于这两个映射互为逆映射，它们是同胚。映射 \(\gamma \mapsto \gamma \delta_0\) 的情形类似。另见 IV, p.~374。
\end{enumerate}

\subsection{庞加莱群胚的函子性}

设 X、Y 是拓扑空间，\(f: X \to Y\) 是连续映射。映射 \(c \mapsto f \circ c\) 是连续映射，记为 \(\Lambda(f)\)，从 \(\Lambda(X) = \mathcal{C}_c(\mathbf{I}; X)\) 到 \(\Lambda(Y) = \mathcal{C}_c(\mathbf{I}; Y)\)（I, p.~132, 引理）。通过取子集，它定义下列连续映射

\[
\Lambda_ {x} (f) \colon \Lambda_ {x} (\mathrm{X}) \to \Lambda_ {f (x)} (\mathrm{Y}),
\]

对 \(x \in X\) 连续，

\[
\Lambda_ {x, y} (f) \colon \Lambda_ {x, y} (\mathrm{X}) \to \Lambda_ {f (x), f (y)} (\mathrm{Y}),
\]

对 \(x, y \in \mathrm{X}\) 连续，并且

\[
\Omega (f) \colon \Omega (\mathrm{X}) \to \Omega (\mathrm{Y}).
\]

对 \(x \in X\)，映射 \(\Lambda_{x,x}(f)\) 也记为 \(\Omega_x(f)\)。取道路连通分支后（III, p.~290），由映射 \(\Lambda_{x,y}(f)\) 得到映射

\[
\varpi_ {x, y} (f) \colon \varpi_ {x, y} (\mathrm{X}) \to \varpi_ {f (x), f (y)} (\mathrm{Y}).
\]

设 \(x, y, z\) 是 X 中的点，\(c \in \Lambda_{x,y}(X)\) 和 \(d \in \Lambda_{y,z}(X)\) 是道路；根据道路拼接的定义，有

\[
f \circ (c * d) = (f \circ c) * (f \circ d).
\]

取严格同伦类后得到关系

\[
\varpi_ {x, z} (f) (\gamma \delta) = (\varpi_ {x, y} (f) (\gamma)) (\varpi_ {y, z} (f) (\delta))
\]

 对所有 \(\gamma \in \varpi_{x,y}(\mathrm{X})\) 及每个 \(\delta \in \varpi_{y,z}(\mathrm{X})\) 都成立。因此，连续映射 \(f\) 以及对 \(\varpi_{x,y}(f)\) 和 \(x\) 的映射 \(y \in \mathrm{X}\)，定义了从群胚 \(\varpi(\mathrm{X})\) 到群胚 \(\varpi(\mathrm{Y})\) 的态射（II, p.~161, 定义 3）。记为 \(\varpi(f)\)，称为由连续映射 \(f\) 诱导的庞加莱群胚态射。特别地，若 \(x\) 是 X 中的一点，则映射 \(\varpi_{x,x}(f)\) 是从群 \(\pi_1(\mathrm{X}, x)\) 到群 \(\pi_1(\mathrm{Y}, f(x))\) 的同态；该同态也记为 \(\pi_1(f, x)\)。

注 1. --- 若在集合 \(\varpi_{x,y}(f)\) 和 \(\varpi_{x,y}(\mathrm{X})\) 上赋予 \(\varpi_{f(x),f(y)}(\mathrm{Y})\) 和 \(\mathcal{C}_c(\mathbf{I};\mathrm{X})\) 上紧开拓扑的商拓扑，则同态 \(\mathcal{C}_c(\mathbf{I};\mathrm{Y})\) 是连续的。

为简化记号，当 \(x\)、\(y \in X\) 且 \(\gamma \in \varpi_{x,y}(X)\) 时，我们有时用 \(f_{*}(\gamma)\) 表示 \(\varpi_{x,y}(f)(\gamma)\) 中的元素 \(\varpi_{f(x),f(y)}(Y)\)。

设 X、Y、Z 是拓扑空间，\(f: X \to Y\) 和 \(g: Y \to Z\) 是连续映射。对 X 中的每条道路 c，有 \((g \circ f) \circ c = g \circ (f \circ c)\)。由此有

\[
\varpi (g \circ f) = \varpi (g) \circ \varpi (f).
\]

设 X 和 Y 是拓扑空间，\(f_{0}\) 和 \(f_{1}\) 是从 X 到 Y 的连续映射，\(\sigma\colon X\times I\to Y\) 是连接 \(f_{0}\) 与 \(f_{1}\) 的同伦。由于从 \((c,t)\mapsto c(t)\) 到 X 的映射 \(\mathcal{C}_{c}(\mathbf{I};\mathrm{X})\times\mathbf{I}\)（III, p.~257, 命题 1）连续，由 \(\mathcal{C}_{c}(\mathbf{I};\mathrm{X})\times\mathbf{I}\times\mathbf{I}\) 给出的从 \((c,t,s)\mapsto\sigma(c(t),s)\) 到 Y 的映射也连续。因此，映射 \(\Sigma\colon(c,s)\mapsto\sigma(c(\cdot),s)\) 是从 \(\mathcal{C}_{c}(\mathbf{I};\mathrm{X})\times\mathbf{I}\) 到 \(\mathcal{C}_{c}(\mathbf{I};\mathrm{Y})\) 的连续映射（上引文）。映射 \(\Sigma\) 是连接映射 \(\Lambda(f_{0})\) 与 \(\Lambda(f_{1})\) 的同伦。将其限制到圈空间，映射 \(\Sigma\) 诱导出同伦 \(\Omega(\mathrm{X})\times\mathbf{I}\to\Omega(\mathrm{Y})\)，连接 \(\Omega(f_{0})\) 与 \(\Omega(f_{1})\)。令 x 为 X 中的一点；假设 \(\sigma\) 是在 x 处的带基点同伦，并置 \(y=f_{0}(x)=f_{1}(x)\)。于是映射 \(\Sigma\) 诱导出从 \(\Omega_{x}(\mathrm{X})\times\mathbf{I}\) 到 \(\Omega_{y}(\mathrm{Y})\) 的连续映射，它是以 \(e_{x}\) 为基点的带基点同伦，连接映射 \(\Omega_{x}(f_{0})\) 与 \(\Omega_{x}(f_{1})\)。

命题 3. --- 设 X 和 Y 是拓扑空间，\(f_{0}\) 和 \(f_{1}\) 是从 X 到 Y 的连续映射，\(\sigma\colon\mathbf{X}\times\mathbf{I}\to\mathbf{Y}\) 是连接 \(f_{0}\) 与 \(f_{1}\) 的同伦。令 x 为 X 中的一点；置 \(y_{0}=f_{0}(x)\)、\(y_{1}=f_{1}(x)\)，并将 \(\delta\in\varpi_{y_{0},y_{1}}(\mathbf{Y})\) 记为由 \(d(t)=\sigma(x,t)\)（\(t\in\mathbf{I}\)）定义的道路 d 的类。对所有 \(\gamma\in\pi_{1}(\mathbf{X},x)\)，有 \((f_{1})_{*}(\gamma)=\delta^{-1}\big((f_{0})_{*}(\gamma)\big)\delta\)。

设 c 是 X 中 x 处的一个圈，\(\gamma\) 是它的严格同伦类。置 \(\gamma_{0} = (f_{0})_{*}(\gamma)\) 和 \(\gamma_{1} = (f_{1})_{*}(\gamma)\)；它们分别是 \(f_{0} \circ c\) 和 \(f_{1} \circ c\) 的严格同伦类。对 \((t, s) \in \mathbf{I} \times \mathbf{I}\)，置 \(\varphi(t, s) = \sigma(c(t), s)\)。对所有 \(t \in I\)，有 \(\varphi(t, 0) = (f_{0} \circ c)(t)\)、\(\varphi(t, 1) = (f_{1} \circ c)(t)\) 且 \(\varphi(0, t) = \varphi(1, t) = d(t)\)。因此关系 \(\gamma_{0}\delta = \delta\gamma_{1}\) 由下述引理得到。

引理 1. --- 设 Y 是一个拓扑空间，\(\varphi\colon\mathbf{I}\times\mathbf{I}\to\mathbf{Y}\) 是连续映射。对 \(t\in\mathbf{I}\)，置 \(c_{0}(t)=\varphi(t,0)\)、\(c_{1}(t)=\varphi(t,1)\)、\(d_{0}(t)=\varphi(0,t)\) 和 \(d_{1}(t)=\varphi(1,t)\)。道路 \(c_{0}*d_{1}\) 与 \(d_{0}*c_{1}\) 严格同伦。

令 c 为在 \(I \times I\) 中通过拼接道路 \(t \mapsto (t,0)\) 和 \(t \mapsto (1,t)\) 得到的道路；令 d 为在 \(I \times I\) 中通过拼接道路 \(t \mapsto (0, t)\) 和 \(t \mapsto (t, 1)\) 得到的道路。道路 c 和 d 具有相同的起点 \(c\)、\(d\)、\((0, 0)\) 和相同的终点 \((1, 1)\)。从 \((t, s) \mapsto (1 - s)c(t) + sd(t)\) 到 \(\mathbf{I} \times \mathbf{I}\) 的映射 \(\mathbf{I} \times \mathbf{I}\) 是连接 c 与 d 的严格同伦。有 \(c\)、\(d\)、\(c_0 * d_1 = \varphi \circ c\) 和 \(d_0 * c_1 = \varphi \circ d\)；因此这两条道路严格同伦。

推论 1. --- 设 X 和 Y 是拓扑空间，x 是 X 中的一点，\(f_0\) 和 \(f_1\) 是从 X 到 Y 的连续映射。若存在在 x 处连接 \(f_0\) 与 \(f_1\) 的带基点同伦，则有 \(\pi_1(f_0, x) = \pi_1(f_1, x)\)。

设 \(\sigma\) 是在 \(x\) 处连接 \(f_0\) 与 \(f_1\) 的带基点同伦。用命题 3 的记号，\(\delta\) 是像为 \(f_0(x) = f_1(x)\) 的常值道路的类。断言由此得到。

推论 2. --- 设 X 和 Y 是拓扑空间，\(f\colon X \to Y\) 是同伦等价。对 X 中的每个点 \(x\)，同态

\[
\pi_ {1} (f, x) \colon \pi_ {1} (\mathrm{X}, x) \to \pi_ {1} (\mathrm{Y}, f (x))
\]

是同构。

设 \(g\)、\(f\) 是从 Y 到 X 的连续映射，是 f 的同伦逆。令 \(x\) 为 X 中的一点。将命题 3 应用于同伦的映射 \(\mathrm{Id}_{\mathrm{X}}\) 和 \(g \circ f: \mathrm{X} \to \mathrm{X}\)，可知映射 \(\pi_1(g \circ f, x)\) 是从群 \(\pi_1(\mathrm{X}, x)\) 到群 \(\pi_1(\mathrm{X}, g \circ f(x))\) 的同构。由于

\[
\pi_ {1} (g \circ f, x) = \pi_ {1} (g, f (x)) \circ \pi_ {1} (f, x),
\]

同态 \(\pi_{1}(f,x)\) 是单射，而同态 \(\pi_{1}(g,f(x))\) 是满射。由于映射 g 也是同伦等价，同态 \(\pi_{1}(g,f(x))\) 是单射；因而它是同构。因此，\(\pi_{1}(f,x)\) 是同构。

例. --- 设 G 是一个拓扑群，e 是它的单位元。对 G 中的每个点 g，左平移和右平移 \(x \mapsto gx\) 及 \(x \mapsto xg\) 都是 G 到自身的同胚（TG, III, p.~2），并将 e 映为 g。由推论 2，它们诱导从 \(\pi_{1}(G, e)\) 到 \(\pi_{1}(G, g)\) 的同构。这些同构不一定相等（IV, p.~459, 习题 1）。

推论 3. --- 设 X 是一个同伦等价于一点的拓扑空间。对 X 中的每个点 x，群 \(\pi_{1}(X,x)\) 退化为单位元。

推论 3 特别适用于 X 是 n 维数值空间 \(R^{n}\) 的情形，更一般地，适用于 X 是 \(R^{n}\) 的一个关于其某个点呈星形的子集的情形（III, p.~234）。

命题 4. --- 设 X 是拓扑空间族 \((\mathrm{X}_{j})_{j\in\mathrm{J}}\) 的积空间。由态射族 \(\varpi(\mathrm{X})\) 定义的从群胚 \(\varpi(\mathrm{X}_{j})\) 到群胚 \(j\in J\)（\((\varpi(\mathrm{pr}_{j}))_{j\in\mathrm{J}}\)）的积的态射是同构。

将此群胚态射记为 \(\varphi\)。它取顶点后诱导的映射是恒等映射 \(X \to \prod_{j} X_{j}\)。设 \(x = (x_{j})\) 和 \(y = (y_{j})\) 是 X 中的两个点。记 \(X\)、\(\varphi_{x,y}\) 为由 \(\varpi_{x,y}(X)\) 诱导的从 \(\prod_{j} \varpi_{x_{j},y_{j}}(X_{j})\) 到 \(\varphi\) 的映射。若对每个 \(j \in J\)，\(c_{j}: I \to X_{j}\) 是连接 \(x_{j}\) 与 \(y_{j}\) 的道路，则映射 \(t \mapsto (c_{j}(t))\) 是 X 中连接 x 与 y 的道路（TG, I, p.~25, 命题 1）。这证明了 \(X\)、\(x\)、\(y\)、\(\varphi_{x,y}\)\(c\)、\(d\)、\(X\)、\(x\)、\(y\)、 是满射。设 c 和 d 是 X 中连接 x 与 y 的两条道路。假设对每个 \(j \in J\)，存在连接 \(\sigma_{j}: I \times I \to X_{j}\) 与 \(\operatorname{pr}_{j} \circ c\) 的严格同伦 \(\operatorname{pr}_{j} \circ d\)。于是，从 \((t,s) \mapsto (\sigma_{j}(t,s))\) 到 X 的映射 \(I \times I\) 是连接 c 与 d 的严格同伦（上引文）。这证明了 \(X\)、\(c\)、\(d\)、\(\varphi_{x,y}\) 是单射。

推论. --- 设 \(x = (x_{j})_{j \in \mathrm{J}}\) 是 X 中的一点。映射

\[
\pi_ {1} (X, x) \to \prod_ {j \in J} \pi_ {1} (X _ {j}, x _ {j})
\]

由映射 \(\pi_1(\mathrm{pr}_j, x_j)\) 诱导，是群同构。

这个同构称为典范同构。下文中，我们常借助此同构将 \(\pi_{1}(\mathrm{X},x)\) 与 \(\prod_{j\in\mathrm{J}}\pi_{1}(\mathrm{X}_{j},x_{j})\) 识别。

注. --- 2) 设 \((\mathrm{X}_{j})_{j\in\mathrm{J}}\) 是一个拓扑空间族。记 X 为积拓扑空间 \(\prod_{j\in J}X_{j}\)，\(x=(x_{j})\) 是 X 中的一点。

对每个 \(i \in \mathrm{J}\)，令 \(u_i: \mathrm{X}_i \to \mathrm{X}\) 为满足下列条件的映射：对 \(z \in \mathrm{X}_i\)，\(\operatorname{pr}_i \circ u_i(z) = z\)，而对所有异于 i 的 \(j \in \mathrm{J}\)，\(i\)、\(\operatorname{pr}_j \circ u_i(z) = x_j\)。映射 \(u_i\) 是连续的，映射 \(\pi_1(u_i, x_i)\) 可与族 \(\pi_1(\mathrm{X}_i, x_i)\) 的积群中因子 \((\pi_1(\mathrm{X}_j, x_j))_{j \in \mathrm{J}}\) 的典范嵌入识别（参见 A, I, p.~45）。

假设集合 J 有限，并且对每个 \(j \in J\)，令 \(\gamma_{j}\) 为 \(\pi_{1}(X_{j}, x_{j})\) 的一个元素。\((\gamma_{j})\) 的元素 \(\pi_{1}(X, x)\) 是圈类 \((u_{j})_{*}(\gamma_{j})\)（\(j \in J\)）的复合，而这些圈类两两可交换。

\begin{enumerate}
\def\labelenumi{\arabic{enumi})}
\setcounter{enumi}{2}
\tightlist
\item
  设 \((\mathrm{X}_j)_{j\in \mathrm{J}}\) 是一个拓扑空间族。记 X 为积拓扑空间 \(\prod_{j\in \mathrm{J}}\mathrm{X}_j\)，\(x = (x_{j})\) 是 X 中的一点。
\end{enumerate}

在集合 \(\pi_1(\mathrm{X},x)\) 和 \(\pi_1(\mathrm{X}_j,x_j)\) 上赋予道路空间 \(\Lambda_x(\mathrm{X})\) 和 \(\Lambda_{x_j}(\mathrm{X}_j)\) 上紧开拓扑的商拓扑。于是同构 \(\pi_1(\mathrm{X},x)\to \prod_{j\in \mathrm{J}}\pi_1(\mathrm{X}_j,x_j)\) 是同胚。它是连续的（III, p.~294, 注 1）。\(\Lambda (\mathrm{X})\) 上的紧开拓扑由形如 \(T(\mathrm{K},\mathrm{U})\) 的子集生成，其中 K 是 I 的紧子集，U 是 X 的开子集。对 \(j\in \mathrm{J}\)，令 \(\mathrm{U}_j\) 为 \(\mathrm{X}_j\) 的开子集，满足 \(\prod_{j\in \mathrm{J}}\mathrm{U}_j\subset \mathrm{U}\)。于是 \((\mathrm{pr}_j)_*(T(\mathrm{K},\mathrm{U}))\) 包含 \(T(\mathrm{K},\mathrm{U}_j)\)。这说明映射 \((\mathrm{pr}_j)_*: \Lambda_x(\mathrm{X})\to \Lambda_{x_j}(\mathrm{X}_j)\) 是开的，映射 \(\pi_1(\mathrm{pr}_j,x_j)\) 也都是开的。由于它们是满射，映射 \(\pi_1(\mathrm{X},x)\to \prod_{j\in \mathrm{J}}\pi_1(\mathrm{X}_j,x_j)\) 是开的（TG, I, p.~34, 命题 8）。它既连续又双射，因而是同胚（TG, I, p.~30, 例 2）。

命题 5. — 设 X 是一个拓扑空间，\((\mathrm{A}_{i})_{i\in\mathrm{I}}\) 是由过滤有序集 I 索引的 X 的子集递增族，并且 X 的每个拟紧子集都包含于某个 \(A_{i}\) 中。典范群胚态射

\[
\rho \colon \varinjlim_ {i \in \mathrm{I}} \varpi (\mathrm{A} _ {i}) \to \varpi (\mathrm{X}),
\]

由 \(A_{i}\) 到 X 的典范嵌入诱导，是同构。若 \(i \leqslant j\)，记 \(\rho_{j,i}\) 为由 \(\varpi(\mathrm{A}_{i}) \to \varpi(\mathrm{A}_{j})\) 到 \(A_{i}\) 的嵌入诱导的群胚态射 \(A_{j}\)。由于 \(\rho_{j,i}\) 取顶点后诱导的映射是嵌入 \(A_{i} \to A_{j}\)，且各 \(A_{i}\) 覆盖 X，\(\rho\) 取顶点后诱导的映射是双射。

设 \(a\) 和 \(b\) 是 X 中的点，\(c\) 是 X 中连接 \(a\) 与 \(b\) 的道路。由于 \(c\) 是紧的，\(\mathbf{I}\) 的像是 X 的拟紧子集（TG, I, p.~62, 定理 2）。因此存在元素 \(i \in \mathrm{I}\)，使得 \(c\) 的像包含于 \(\mathrm{A}_i\) 中。于是，\(\rho\) 取箭集后诱导的映射是满射。

令 \(i \in \mathrm{I}\)，令 \(a\) 和 \(b\) 为 \(\mathrm{A}_i\) 中的点，令 \(c\) 和 \(c'\) 为 \(a\) 中连接 \(b\) 与 \(\mathrm{A}_i\) 的道路；令 \(h\) 为 \(c\) 中连接 \(c'\) 与 \(\mathrm{X}\)、\(\mathbf{I} \times \mathbf{I}\) 的严格同伦。由于 \(h(\mathbf{I} \times \mathbf{I})\) 是紧的，\(\mathrm{X}\) 是 \(i \in \mathrm{I}\) 的拟紧子集（上引文），并且存在元素 \(h\)，使得 h 的像包含于 \(\mathrm{A}_i\)、\(c\) 中。道路 c 与 \(c'\) 在 \(A_{i}\) 中严格同伦；从而，道路类 {[}c{]} 与 \([c']\) 在 \(\varinjlim\varpi(A_{i})\) 中的像相同。因此，\(\rho\) 是单射。

推论. --- 设 a 是 X 中的一点，J 是由满足 \(i \in I\) 的 \(a \in A_{i}\) 组成的集合。典范同态

\[
\varinjlim_ {i \in J} \pi_ {1} (A _ {i}, a) \to \pi_ {1} (X, a)
\]

是双射。

注 4. --- 当 \((\mathrm{A}_{i})_{i\in\mathrm{I}}\) 是由过滤有序集 I 索引的 X 的子集递增族，且各 \(A_{i}\) 的内部覆盖 X 时，本命题及其推论尤其适用。

\subsection{自由同伦的圈}

定义 4. --- 设 X 是一个拓扑空间，c 和 c' 是 X 中的两个圈。若同伦 σ 连接 c 与 c'，且对所有 s ∈ I 都有 σ(0, s) = σ(1, s)，则称其为连接 c 与 c' 的自由同伦。如果存在连接 c 与 c' 的自由同伦，就称 c 与 c' 自由同伦。

连接 \(c\) 与 \(c'\) 的自由同伦对应于 X 的圈空间 \(\Omega(\mathrm{X})\) 中连接 \(c\) 与 \(c'\) 的道路。因此，关系“ \(c\) 与 \(c'\) 自由同伦”是 \(\Omega(\mathrm{X})\) 上的等价关系，其等价类是 \(\Omega(\mathrm{X})\) 的道路连通分支。

注. --- 设 \(\varphi\) 是从 \(\mathbf{R}\) 到 \(\mathbf{T} = \mathbf{R}/\mathbf{Z}\) 的典范映射（TG, V, p.~2）。映射 \(f \mapsto f \circ \varphi|\mathbf{I}\) 是从 \(\mathcal{C}_c(\mathbf{T};\mathbf{X})\) 到 \(\Omega(\mathbf{X})\) 的同胚，从而取道路连通分支后，得到从集合 {[}T; X{]}（III, p.~230）到 X 中的自由同伦圈类集合的双射。

命题 6. --- 设 X 是一个道路连通的拓扑空间，x 是 X 中的一点。

\begin{enumerate}
\def\labelenumi{\alph{enumi})}
\item
  X 中的每个圈都与 x 处的一个圈自由同伦。更确切地说，若 c 是 y 处的圈，d 是起点为 y、终点为 x 的道路，则 c 与 x 处的圈 \((\overline{d} * c) * d\) 自由同伦。
\item
  X 中 x 处的两个圈自由同伦，当且仅当它们的严格同伦类在群 \(\pi_{1}(X,x)\) 中共轭。
\end{enumerate}

我们来证明 \(a\)。对每个 \(s \in [0,1]\)，记 \(d_s\) 为 X 中由 \(d_s(t) = d(st)\)（\(t \in \mathbf{I}\)）定义的道路；其起点为 \(y\)。由于映射 \((s,t) \mapsto d(st)\) 连续，从 \(s \mapsto d_s\) 到 \(\mathbf{I}\) 的映射 \(\mathcal{C}_c(\mathbf{I};\mathbf{X})\) 是连续的（III, p.~257, 命题 1）。于是映射 \(s \mapsto (\overline{d_s} * c) * d_s\) 是 \(\Omega(\mathbf{X})\) 中的一条道路（III, p.~257, 命题 2），连接 \((e_y * c) * e_y\) 与 \((\overline{d} * c) * d\)，故得 \(a\)。

设 c 和 \(c'\) 是 X 中 x 处的两个圈。若它们的严格同伦类在 \(\pi_{1}(X,x)\) 中共轭，则存在 x 处的圈 d，使得 \(c'\) 与圈 \((\overline{d} * c) * d\) 严格同伦。由 a)，c 与 \(c'\) 自由同伦。反之，假设存在连接 c 与 \(\varphi\) 的自由同伦 \(c'\)。置 \(d(t) = \varphi(0,t)\)，则也有 \(d(t) = \varphi(1,t)\)，且 d 是 x 处的圈。由 III, p.~295 的引理 1，圈 \(c * d\) 与 \(d * c'\) 严格同伦。因此，c 与 \(c'\) 的严格同伦类在 \(\pi_{1}(X,x)\) 中共轭。

系理. --- 设 X 是一个道路连通的拓扑空间，x 是 X 中的一点。命题 6 使我们能够定义从 X 中的自由同伦圈类集合到 \(\pi_{1}(X,x)\) 中的共轭类集合的典范双射。

\section{同伦与覆叠}

\subsection{同伦与覆叠}

命题 1. --- 设 B 是一个拓扑空间，E 是 B 的一个覆叠。设 B' 是一个拓扑空间，\(f_0\) 和 \(f_1\) 是从 B' 到 B 的连续映射。若映射 \(f_0\) 和 \(f_1\) 同伦，则 B' 的覆叠 \(f_0^*(\mathrm{E})\) 和 \(f_1^*(\mathrm{E})\) 同构。

设 \(\sigma\colon\mathrm{B}^{\prime}\times\mathbf{I}\to\mathrm{B}\) 是连接 \(f_{0}\) 与 \(f_{1}\) 的同伦。记 \(i_{0}\) 和 \(i_{1}\) 为从 \(x\mapsto(x,0)\) 到 \(x\mapsto(x,1)\) 的映射 \(\mathrm{B}^{\prime}\) 和 \(\mathrm{B}^{\prime}\times\mathbf{I}\)。若 \(t\in\{0,1\}\)，则有 \(f_{t}=\sigma\circ i_{t}\)，因而 \(f_{t}^{*}(\mathrm{E})\) 的覆叠 \(i_{t}^{*}(\sigma^{*}(\mathrm{E}))\) 和 \(\mathrm{B}^{\prime}\) 同构（I, p.~15）。由于拓扑空间 \(\mathbf{I}\) 局部连通且单连通（I, p.~127, 推论），覆叠 \(i_{0}^{*}(\sigma^{*}(\mathrm{E}))\) 的 \(i_{1}^{*}(\sigma^{*}(\mathrm{E}))\) 和 \(B'\) 同构（I, p.~130, 命题 8 的推论 1）。

推论. --- 同胚于单连通拓扑空间的拓扑空间是单连通的。

设 B 和 B' 是拓扑空间，\(f: B \rightarrow B'\) 是同伦等价，\(g: B' \rightarrow B\) 是 f 的同伦逆连续映射。设 E 是 B 的一个覆叠。由于 \(g \circ f\) 同伦于 B 的恒等映射，覆叠 E 与覆叠 \(f^{*}(g^{*}(\mathrm{E}))\) 同构（命题 1）。若空间 B' 是单连通的，则覆叠 \(g^{*}(\mathrm{E})\) 可平凡化。因此覆叠 E 可平凡化。这证明空间 B 是单连通的。

注. --- 设 G 是一个离散拓扑群。用本命题的记号，假设 E 是群 G 的主覆叠。于是覆叠 \(f_{0}^{*}(E)\) 和 \(f_{1}^{*}(E)\) 是同构的主覆叠（I, p.~131, 注）。

命题 2（同伦的提升）. --- 设 B 是一个拓扑空间，E 是 B 的一个覆叠。设 B' 是一个拓扑空间，\(f_0\) 和 \(f_1\) 是从 B' 到 B 的连续映射。设 \(\sigma: \mathrm{B}' \times \mathbf{I} \to \mathrm{B}\) 是连接 \(f_0\) 与 \(f_1\) 的同伦。设 \(g_0: \mathrm{B}' \to \mathrm{E}\) 是 \(f_0\) 到 E 的连续提升。存在唯一的连续映射 \(\widetilde{\sigma}: \mathrm{B}' \times \mathbf{I} \to \mathrm{E}\) 提升 \(\sigma\)，且它是以 \(g_0\) 为起点的同伦。它的终点 \(g_1: \mathrm{B}' \to \mathrm{E}\) 是 \(f_1\) 到 E 的连续提升。

这由 I, p.~130 的命题 8 的推论 3 应用于单连通且局部连通的拓扑空间 I 而得。

\subsection{道路的提升}

命题 3. --- 设 B 是一个拓扑空间，\(p\colon E \to B\) 是一个覆叠。记 \(\widetilde{p}\colon \mathcal{C}_{c}(\mathbf{I};E) \to \mathcal{C}_{c}(\mathbf{I};B)\) 为映射 \(c \mapsto p \circ c\)。记 \(o_{\mathrm{E}}\colon \mathcal{C}_{c}(\mathbf{I};E) \to E\)（分别记 \(o_{\mathrm{B}}\colon \mathcal{C}_{c}(\mathbf{I};B) \to B\)）为由 \(c \mapsto c(0)\) 定义的映射。则图

\[
\begin{array}{c} \mathcal {C} _ {c} (\mathbf {I}; \mathrm{E}) \xrightarrow {o _ {\mathrm{E}}} \mathrm{E} \\ \Big \downarrow \widetilde {p} \\ \mathcal {C} _ {c} (\mathbf {I}; \mathrm{B}) \xrightarrow {o _ {\mathrm{B}}} \mathrm{B} \end{array}
\]

是一个笛卡尔方块。

拓扑空间 I 局部连通、局部紧且单连通。因此，将 I, p.~131 的命题 9 应用于 T = I、t = 0，即得本命题。

推论 1. --- 映射 \(\widetilde{p}\colon\mathcal{C}_{c}(\mathbf{I};\mathrm{E})\to\mathcal{C}_{c}(\mathbf{I};\mathrm{B})\) 是一个覆叠。

这由 I, p.~71 的命题 1 的推论 2 得到。

推论 2（道路的提升）. --- 设 \(p\) : E → B 是一个覆叠，x 是 E 中的一点，a = p(x)。映射 c ↦ p ∘ c 是从 E 中起点为 x 的道路空间 \(\Lambda_x\) (E) 到 B 中起点为 a 的道路空间 \(\Lambda_a\) (B) 的同胚。

用命题 3 的记号，有 \(\Lambda_{a}(\mathrm{B}) = o_{\mathrm{B}}^{-1}(a)\) 和 \(\Lambda_x(\mathrm{E}) = o_{\mathrm{E}}^{-1}(x)\)。推论于是由 I, p.~10 的命题 4 的推论得到。

如果连续映射 \(p \colon \mathrm{E} \to \mathrm{B}\) 满足下述道路提升性质，我们就这样称呼它：对每条道路 \(c \colon \mathbf{I} \to \mathbf{B}\) 以及 E 中提升 \(x\) 的每个点 \(c(0)\)，存在一条 E 中起点为 \(c'\) 且提升 c 的道路 \(x\)。若 \(c\)、\(p\) 是平展且分离的，这样的道路便是唯一的（I, p.~34, 命题 11 的推论 1）。设 E 是一个覆叠，\(p\) 是其投影；映射 \(p\) 是平展、分离的，并且具有道路提升性质（推论 2）。部分逆命题见 III, p.~315 的推论。

命题 4. --- 设 \(p \colon \mathrm{E} \to \mathrm{B}\) 是满足道路提升性质的平展分离映射。设 \(a\) 和 \(b\) 是 \(\mathrm{B}\) 中的点，\(x\) 是 \(\mathrm{E}\) 中满足 \(p(x) = a\)、\(c_0\)、\(c_1\) 的点。两条在 \(a\) 中连接 \(b\) 与 \(\mathrm{B}\) 的严格同伦道路，其分别以 x 为起点在 \(x\)、\(\mathrm{E}\)、\(c_0\)、\(c_1\) 中的提升具有相同的终点，并且严格同伦。

设 \(\sigma\colon\mathbf{I}\times\mathbf{I}\to\mathbf{B}\) 是连接 \(c_{0}\) 与 \(c_{1}\) 的严格同伦。对 \(s\in\mathbf{I}\)，令 \(c_{s}^{\prime}\) 为 E 中以 x 为起点、提升道路 \(x\)、\(t\mapsto\sigma(t,s)\) 的唯一道路。对 \((t,s)\in\mathbf{I}\times\mathbf{I}\)，置 \(\sigma^{\prime}(t,s)=c_{s}^{\prime}(t)\)；有 \(p\circ\sigma^{\prime}=\sigma\)。映射 \(\sigma^{\prime}\) 在 \(\{0\}\times\mathbf{I}\) 上为常值；按构造，它在每个 \(\mathbf{I}\times\{s\}\) 上连续。因此由 I, p.~37, 定理 1 的推论 1，它是连续的。特别地，从 \(s\in\mathbf{I}\)、\(\mathbf{I}\) 到 \(\mathrm{E}_{b}\) 的映射 \(s\mapsto\sigma^{\prime}(1,s)\)、\(b\) 是提升像为 b 的常值道路的连续映射；因而它是常值的。这证明映射 \(\sigma^{\prime}\) 是连接 \(c_{0}^{\prime}\) 与 \(c_{1}^{\prime}\) 的严格同伦。

推论 1. --- 设 B 是一个拓扑空间，E 是 B 的一个覆叠；记 p 为其投影。对 E 中的每一对 \((x,y)\)，映射 \(\varpi_{x,y}(p)\colon\varpi_{x,y}(\mathrm{E})\to\varpi_{p(x),p(y)}(\mathrm{B})\) 是单射。特别地，对 E 中的每个点 x，同态 \(\pi_{1}(p,x)\colon\pi_{1}(\mathrm{E},x)\to\pi_{1}(\mathrm{B},p(x))\) 是单射。

推论 2. --- 设 B 是一个拓扑空间，E 是 B 的一个覆叠；记 p 为其投影。令 x 为 E 中的一点，置 \(b = p(x)\)。对于 B 中 a 处的圈 c，其在 \(\pi_{1}(B, b)\) 中的类属于同态 \(\pi_{1}(p, x)\) 的像子群，当且仅当以 x 为起点提升 c 的道路 c' 是 E 中 x 处的圈。

这个条件显然是充分的。反之，设 \(c_{1}^{\prime}\) 是 E 中 x 处的一个圈。置 \(c_{1} = p \circ c_{1}^{\prime}\)，并假设圈 c 与 \(c_{1}\) 严格同伦。由命题 4，道路 \(c^{\prime}\) 与道路 \(c_{1}^{\prime}\) 具有相同的终点。因此它是 x 处的圈。

推论 3. --- 设 B 是一个拓扑空间，\((\mathrm{E}_{1}, p_{1})\) 和 \((\mathrm{E}_{2}, p_{2})\) 是 B 的覆叠。记 \((\mathrm{E}, p)\) 为它们的 B-积空间，\(x = (x_{1}, x_{2})\) 是 E 中的一点。同态 \(\pi_{1}(p, x)\) 的像等于同态 \(\pi_{1}(p_{1}, x_{1})\) 和 \(\pi_{1}(p_{2}, x_{2})\) 的像的交。

设 \(c\) 是 B 中以 \(p_{1}(x_{1}) = p_{2}(x_{2})\) 为基点的圈，并且对 \(i \in \{1, 2\}\)，令 \(c_{i}^{\prime}\) 为 \(x_i\) 中以 \(E_i\) 为起点、提升 c 的道路。B-空间 \(c\)、\(E = E_{1} \times_{B} E_{2}\) 是 B 的一个覆叠（I, p.~72, 命题 2 的推论 3），而道路 \(c^{\prime}: t \mapsto (c_{1}^{\prime}(t), c_{2}^{\prime}(t))\) 是 E 中以 x 为起点提升 c 的唯一道路。\(x\)、\(c\)、\(c^{\prime}\) 是圈的充要条件是 \(c_{1}^{\prime}\) 和 \(c_{2}^{\prime}\) 都是圈。因此推论 3 由推论 2 得到。

\subsection{覆叠空间中庞加莱群胚的作用}

设 B 是一个拓扑空间，\(p \colon \mathrm{E} \to \mathrm{B}\) 是满足道路提升性质的分离平展映射（III, p.~302）；如果映射 \(p\) 使 E 成为 B 的覆叠，情形就是如此（上引文）。设 \(b\) 和 \(b'\) 是 B 中的点，\(c \in \Lambda_{b,b'}(\mathrm{B})\) 是 B 中连接 \(b\) 与 \(b'\) 的道路。对纤维 \(x\)、\(\mathrm{E}_b\) 中的每个点 \(x \cdot c\)，记 E 中以 x 为起点且提升 c 的道路的终点为 \(x\)。点 \(c\) 属于纤维 \(x \cdot c\)，并且只依赖于道路 c 的类 \(E_{b'}\)（III, p.~302, 命题 4）；因此将 \(\gamma \in \varpi_{b,b'}(\mathrm{B})\)、\(x \cdot \gamma\) 记为 \(x \cdot c\)。

若 c 是 b 处的常值道路，则有 \(c\)、\(b\)、\(x \cdot \gamma = x\)、\(x\)、\(c\)，因为以 x 为起点的常值道路提升 c。

设 \(b''\) 是 B 中的另一个点，\(c' \in \Lambda_{b',b''}(B)\)。对 \(x\) 中的每个点 \(E_b\)，有：

\[
(x \cdot c) \cdot c ^ {\prime} = x \cdot (c * c ^ {\prime}).\tag{1}
\]

事实上，记 \(\widetilde{c}\) 为以 x 为起点的 c 的提升，记 \(x\)、\(c\)、\(\widetilde{c}'\)、\(c'\) 为以 \(x \cdot c = \widetilde{c}(1)\) 为起点的 c' 的提升。道路 \(\widetilde{c}\) 与 \(\widetilde{c}'\) 可拼接，且 \(\widetilde{c} * \widetilde{c}'\) 是以 x 为起点提升 \(x\)、\(c * c'\) 的道路。

记 \(\varphi_{b,b'}\colon \varpi_{b,b'}(\mathrm{B})\to \mathcal{F}(\mathrm{E}_b;\mathrm{E}_{b'})\) 为满足下列条件的映射：对 \(\gamma \in \varpi_{b,b'}(\mathrm{B})\) 和 \(x\in \mathrm{E}_b\)，有

\[
\varphi_ {b, b ^ {\prime}} (\gamma) (x) = x \cdot \gamma .
\]

对每个 \(\gamma \in \varpi_{b,b'}(B)\) 和每个 \(\gamma' \in \varpi_{b',b''}(B)\)，由关系 (1) 可得

\[
\varphi_ {b, b ^ {\prime \prime}} (\gamma \gamma^ {\prime}) = \varphi_ {b ^ {\prime}, b ^ {\prime \prime}} (\gamma^ {\prime}) \circ \varphi_ {b, b ^ {\prime}} (\gamma).\tag{2}
\]

族 \(\varphi = (\varphi_{b,b'})_{(b,b')\in\mathrm{B}\times\mathrm{B}}\) 因而关于映射 \(\varpi(\mathrm{B})\) 在集合 E 上定义了群胚 \(p\colon\mathrm{E}\to\mathrm{B}\) 的右作用法则（II, p.~167）。称其为与映射 \(\varpi(\mathrm{B})\) 相伴的群胚 \(p\colon\mathrm{E}\to\mathrm{B}\) 的典范作用。映射 \(\varphi_{b,b}\colon\pi_{1}(\mathrm{B},b)\to\mathcal{F}(\mathrm{E}_{b};\mathrm{E}_{b})\) 定义群 \(\pi_{1}(\mathrm{B},b)\) 在纤维 \(\mathrm{E}_{b}\) 上的右作用。称此作用为 \(\pi_{1}(\mathrm{B},b)\) 在纤维 \(\mathrm{E}_{b}\) 上的典范作用。

注. --- 设 \(b\) 和 \(b'\) 是 B 中的两个点，\(x\) 是纤维 \(\mathrm{E}_b\) 中的一点，\(c \in \Lambda_{b,b'}(\mathrm{B})\) 是 B 中的一条道路，\(\widetilde{c}\)、\(x\)、\(c\) 是 E 中以 x 为起点且提升 c 的道路。对所有 \(s \in \mathbf{I}\)，记 \(c_s\) 为道路 \(t \mapsto c(st)\)。E 中以 x 为起点提升 \(x\)、\(c_s\) 的道路是 \(t \mapsto \widetilde{c}(st)\)，其终点是点 \(\widetilde{c}(s)\)。因此，对所有 \(\widetilde{c}(s) = x \cdot c_s\)，有 \(s \in \mathbf{I}\)。

设 B' 是一个拓扑空间，\(p'\colon E'\to B'\) 是满足道路提升性质的平展分离映射。设 \(f\colon B\to B'\) 和 \(g\colon E\to E'\) 是连续映射，满足 \(p'\circ g=f\circ p\)。对 \(b, b'\in B, \gamma\in\varpi_{b,b'}(B)\) 和 \(x\in E_b\)，有：

\[
g (x \cdot \gamma) = g (x) \cdot f _ {*} (\gamma).\tag{3}
\]

事实上，令 c 为 B 中严格同伦类为 \(\gamma\) 的道路；记 \(\widetilde{c}\) 为 E 中以 x 为起点且提升 c 的道路。于是道路 \(g \circ \widetilde{c}\) 是 E' 中以 \(g(x)\) 为起点提升 \(f \circ c\) 的道路。

特别地，当 \(B = B'\) 且 \(f = \operatorname{Id}_{B}\) 时，有

\[
g (x \cdot \gamma) = g (x) \cdot \gamma .\tag{4}
\]

设 \(p \colon \mathrm{E} \to \mathrm{B}\) 和 \(p' \colon \mathrm{E}' \to \mathrm{B}\) 是满足道路提升性质的平展分离映射。设 b 是 B 中的一点。若 \(b\)、\(f \colon \mathrm{E} \to \mathrm{E}'\) 是 B-态射，则映射 \(f_b \colon \mathrm{E}_b \to \mathrm{E}_b'\) 是一个 \(\pi_1(\mathrm{B}, b)\)-态射。

定理 1. --- 设 B 是一个拓扑空间， 是 B 的一个覆叠，b 是 B 中的一点，x 是纤维 \(E_{b}\) 中的一点。

\begin{enumerate}
\def\labelenumi{\alph{enumi})}
\item
  x 在群 \(\pi_{1}(\mathrm{B},b)\) 的典范作用下的轨道等于 \(E_{b}\) 与 x 在 E 中的道路连通分支的交。特别地，若空间 E 是道路连通的，则此作用是传递的。
\item
  x 的固定子（A, I, p.~51）是 \(p_{*}(\pi_{1}(E,x))\) 的子群 \(\pi_{1}(B,b)\)。
\item
  对每个 \(\gamma \in \pi_1(\mathrm{B},b)\)，有 \(p_*(\pi_1(\mathrm{E},x)) = \operatorname{Int}(\gamma)(p_*(\pi_1(\mathrm{E},x \cdot \gamma)))\)。
\item
  如果空间 B 是连通的且覆叠 E 是伽罗瓦覆叠，则子群 \(p_{*}(\pi_{1}(\mathrm{E},x))\) 在 \(\pi_{1}(\mathrm{B},b)\) 中是正规子群。
\end{enumerate}

断言 a) 是直接的。断言 b) 由命题 4 的推论 2（III, p.~303）得到。断言 c) 由 b) 以及 A, I, p.~52 的命题 2 得到。最后，假设 B 是连通的，且 E 是 B 的伽罗瓦覆叠。对每个 \(\gamma \in \pi_1(B,b)\)，存在 E 的一个 B-自同构 \(g\)，使得 \(g(x) = x \cdot \gamma\)（I, p.~102, 定理 2, c)）。有 \(p = p \circ g\)，从而 \(p_*(\pi_1(E,x)) = p_*(\pi_1(E,x \cdot \gamma))\)。因此断言 d) 由 c) 得到。
